\cleardoublepage
\chapter{Future outlook}
\phantomsection
%\addcontentsline{toc}{chapter}{Future outlook}

We end by offering an overview on how the author expects the work started in this thesis to be continued. In the first section we will explore how the interval category $\Int$ can be generalized beyond dimension $n=1$. The goal is articulating a category $\cube_n$, which when localized along a certain class of maps produces a category equivalent to $\btheta_n$. Using this conjectured localization we define $(\oo,n)$-categories using $\cube_n$, by introducing a generalization of the interval Segal condition \Cref{ch2:: def: interval Segal condition}. Lastly, we put forth a set of submersivity conditions intended to generalize \Cref{ch3:: def: bord_d} and \Cref{ch5:: def: bord_d strat}, and we conjecture that, together with a concordance-resolution step, it produces a model for the $(\oo,n)$-category of fully-extended (stratified) bordisms.

In the last section, we conjecture the existence of a mechanism to produce stratified $d$-dimensional (non-extended) topological quantum field theories using \emph{stratified factorization homology} — a variant of ordinary factorization homology, first introduced in \cite{LurieHA} as topological chiral homology, and later developed in \cite{AFT1}.\footnote{As of this writing, many authors refer to this also as $\alpha$-factorization homology.} The main goal is to bring forth the first touchstone in the attempt to give a stratified generalization of Scheimbauer's thesis \cite{Sch14}, where the author builds a fully-extended topological quantum field theory using $\alpha$-factorization homology. 

\section{Extending intervals to cubes}

The main idea driving this section is to extend the interval Segal model of \((\infty,1)\)-categories to higher dimensions.  For \((\infty,1)\)-categories we used the category \(\Int\), of finite strings of closed intervals as a Segal-type indexing category playing the role of \(\bDelta\).  In higher dimensions, a natural globular replacement for \(\bDelta\) is Joyal–Berger’s category \(\btheta_n\); Rezk’s \(\Theta_n\)-spaces give a well–known model for
\((\infty,n)\)-categories.

Following an idea learned from unpublished work of Manuel Araújo and Pedro Boavida de Brito, we introduce an \(n\)-dimensional cubical indexing category \(\cube_n\) which plays for \(\btheta_n\) the same role that \(\Int\) plays for \(\bDelta\). Objects of \(\cube_n\) are \(n\)-cubes equipped with systems of time coordinates, given by decompositions of each edge into subintervals; morphisms are generated by face maps, degeneracies and refinements of these decompositions. Crucially, \(\cube_n\) is not simply the product \(\Int \times \cdots \times \Int\), since we allow different faces of a cube to carry independent decompositions.


\subsection{\texorpdfstring{$\bTheta_n$}{Theta_n}-spaces}

Before we define $\cube_n$, we start by reviewing $\bTheta_n$, and the notion of Rezk--Segal $\btheta$-space as a model for $(\oo,n)$-categories.

\subsubsection{The category \texorpdfstring{$\bTheta_n$}{Theta_n}}

The literature is inundated with differing models for $(\oo,n)$-categories. This is made especially salient in the wonderful exposition of Chang--Lauda \cite{CL04} and the comparison works of Bergner--Rezk \cite{BR13,BR14}, using iterated wreath products. The wreath construction can be defined in varying degrees of generality, we will follow the variant first described in the work of Berger \cite{BER05} and Bergner--Rezk.

First, we re-introduce Segal's category of finite sets.

\begin{warning}
We will now \emph{re-define} Segal's category and we will add a remark afterwards about the equivalence with the previously defined one.
\end{warning}

\begin{definition}
    Let $\Gamma$ be the category with the following structure:\
    \begin{itemize}
        \item \textbf{Objects}: Finite sets $X,Y,\ldots$.
        \item\textbf{Morphisms}: A morphism from $X$ to $Y$ is given by a function 
        \begin{equation}
            \theta \maps X \longra \cP(Y)
        \end{equation} 
        where $\cP(Y)$ is the set of all subsets of $Y$, with the additional property that $\theta(x_1) \cap \theta(x_2) = \emptyset$ are disjoint for $x_1\neq x_2$. 
        \item \textbf{Composition}: Given a pair of composable morphism $X\xra{\theta} Y \xra{\mu} Z$, we define the composite $\mu \circ \theta$ by sending
        \begin{equation}
            x \longmapsto  \bigcup_{y \in \theta(x)} \mu(y).
        \end{equation} 
    \end{itemize}
\end{definition}
\begin{proposition}
There is a functor, which we will henceforth denote as the \textbf{Segal functor},
\begin{equation}
\cT \maps \bDelta \longra \Gamma    
\end{equation}
which maps $[n]$ to the underlying finite set $\{0,1\ldots,n\}$, and maps a morphism $\alpha\maps [n] \longra [m]$ to the following map in $\Gamma$:
\begin{equation}
    i \longmapsto \{j \in [m]\mkern1mu : \, \alpha(i-1) < j \leq \alpha(i)\}.
\end{equation}
\end{proposition}
\begin{proof}
    The functor is well-defined since $\cT(f)(i)$ must necessarily be disjoint from $\cT(f)(j)$ for $i\neq j$. Moreover, functoriality is direct from the fact that the composite of order-preserving maps is still order-preserving.
\end{proof}
Having defined Segal's category of finite sets, we can now turn to the definition of \emph{wreath product}. 

\begin{definition}\label{ch6:: def: wreath_product}
    Let $\cD$ be an arbitrary category. We define $\Gamma \wr \cD$, which we will call the \textbf{wreath product}, to be the following category:
    \begin{itemize}
        \item\textbf{Objects}: Objects are \emph{symbols} $X(d_{x})_{x \in X}$, where $X$ is a finite set and $(d_{x})_{x\in X}$ is an element of $\mathsf{Obj}(\cD)^{X}$, i.e., an $X$-indexed sequence of objects of $\cD$.
        \item\textbf{Morphisms}: A morphism 
        \begin{equation}
            \theta \maps X(d_{x})_{x \in X} \longra Y(e_{y})_{y \in Y}  
        \end{equation}
        consists of tuples $(\theta_{\Gamma},\{\theta_{xy}\})$; these are a morphism $\theta_\Gamma \maps X \longra Y$ in $\Gamma$, together with a collection of morphisms 
        \begin{equation}
            \theta_{xy} \maps d_x \longra e_y
        \end{equation}
        in $\mathsf{Hom}(\cF)_{\cD}(d_x,e_y)$ whenever $y\in \theta_{\Gamma}(x)$.
        \item \textbf{Composition:} Given a pair of morphisms $\theta \maps X(d_x)_{x \in X} \longra Y(e_y)_{y\in Y}, \phi \maps Y(e_y)_{y\in Y} \longra Z(c_z)_{z\in Z}$; we define its composite morphism to be:
        \begin{enumerate}
            \item Its $\Gamma$ component is defined as 
            \begin{equation}
                (\phi \circ \theta)_{\Gamma} = \phi_\Gamma \circ \theta_{\Gamma}.
            \end{equation}
            \item The collection of morphisms 
            \begin{equation}
                (\phi \circ \theta)_{xz} \maps d_x \longra c_{z}
            \end{equation}
            is defined as containing all composites $\phi_{yz}\circ \theta_{xy}$, subject to $y \in \theta_\Gamma (x)$ and $z \in \phi_\Gamma(y)$.
        \end{enumerate}
    \end{itemize} 
\end{definition}

There is manifestly a forgetful functor $\Gamma \wr \cD \longra \Gamma$, which sends the symbol $X(d_x)_{x\in X}$ to the finite set $X$, and sends a morphism $\theta$ to the underlying map $\theta_\Gamma\colon X \longra Y$ in $\Gamma$. 

\begin{definition}\label{ch6:: def: wreath_prod_general}
    Let $\cF \longra \Gamma$ be an object in $\Cat/\Gamma$. Let the functorial construction:
    \begin{align}
    (-)\wr(-) \maps (\Cat/\Gamma) \times \Cat &\longra \Cat \\
    (\cF \longra \Gamma, \cD) &\longmapsto \cF \times_{\Gamma} (\Gamma \wr \cD) 
    \end{align} 
    which we will denote as $\cF \wr \cD$, and call it the \textbf{wreath product} of $\cF$ with $\cD$.
\end{definition}

Using the Segal functor $\cT \maps \bDelta \longra \Gamma$, we are now allowed to define the wreath product $\bDelta \wr \cF$ with any category $\cF$. Despite the redundancy in exposition, let us carefully unpack its structure, as it is the keystone to building $\bTheta_n$.

\begin{definition}
    Let $\cF$ be an ordinary 1-category. Let $\bTheta \wr \cF$ be the category with the following structure:
    \begin{itemize}
        \item \textbf{Objects}: Tuples $([m],c_1,\ldots,c_m)$, where $[m]\in \bDelta$ and $c_i \in \cF$.
        \item \textbf{Morphisms}: The set of morphisms $[m](c_1,\ldots,c_m) \longra [n](d_1,\ldots,d_n)$ consists of tuples $(\alpha,\{f_{ij}\})$ such that:
        \begin{enumerate}
            \item $\alpha \maps [m] \longra [n]$ is a morphism in $\bDelta$.
            \item For each pair $(i,j)\in [m]\times [n]$ where $\alpha(i-1) < j \leq \alpha(i)$, a morphism 
            \begin{equation}
                f_{ij} \maps c_i \longra d_j
            \end{equation}
            in $\cF$. 
        \end{enumerate}
        \item \textbf{Composition}: Given a pair of composable morphisms:
        \begin{equation}
            ([m],c_1,\ldots,c_m) \xra{(\alpha,\{f_{ij}\})} ([n],d_1,\ldots,d_n) \xra{(\beta,\{g_{jk}\})}
            ([p],e_1,\ldots,e_p)
        \end{equation}
        we define its composite $(\beta\circ \alpha, \{g \circ f\}_{ik})$ as:
        \begin{enumerate}
            \item $(\beta \circ \alpha) \maps [m] \longra [p]$.
            \item For each triple $(i,j,k) \in [m] \times [n] \times [p]$ where:
            \begin{enumerate}
                \item $\alpha(i-1) < j \leq \alpha(i)$, and
                \item $\beta(j-1) < k \leq \beta(j)$,
            \end{enumerate}
            a morphism $(g \circ f)_{ij}$ defined as the composite $g_{jk}\circ f_{ij}$. 
        \end{enumerate}
    \end{itemize} 
\end{definition}
\begin{observation}
There is a useful diagrammatic way of thinking about objects in $\bDelta \wr \cF$, where we depict the tuples as labelled finite directed graphs, and the morphisms as shown below: 
\begin{equation}
    \begin{tikzcd}[ampersand replacement=\&]
        \& 0 \&\& 1 \&\&\&\& 0 \&\& 1 \&\& 2 \&\& 3 \\
        \\
        0 \&\& 1 \&\& 2 \&\& 0 \&\& 1 \&\& 2 \&\& 3
        \arrow[""{name=0, anchor=center, inner sep=0}, "{c_1}"{description}, from=1-2, to=1-4]
        \arrow[between={0.1}{0.8}, maps to, from=1-2, to=3-1]
        \arrow[between={0.1}{0.8}, maps to, from=1-4, to=3-5]
        \arrow[""{name=1, anchor=center, inner sep=0}, "{c_1}"{description}, from=1-8, to=1-10]
        \arrow[between={0.1}{0.9}, maps to, from=1-8, to=3-7]
        \arrow["{c_2}"{description}, from=1-10, to=1-12]
        \arrow[between={0.1}{0.8}, maps to, from=1-10, to=3-11]
        \arrow[""{name=2, anchor=center, inner sep=0}, "{c_3}"{description}, from=1-12, to=1-14]
        \arrow[between={0.1}{0.8}, maps to, from=1-12, to=3-11]
        \arrow[between={0.1}{0.8}, maps to, from=1-14, to=3-13]
        \arrow[""{name=3, anchor=center, inner sep=0}, "{d_1}"{description}, from=3-1, to=3-3]
        \arrow[""{name=4, anchor=center, inner sep=0}, "{d_2}"{description}, from=3-3, to=3-5]
        \arrow[""{name=5, anchor=center, inner sep=0}, "{d_1}"{description}, from=3-7, to=3-9]
        \arrow[""{name=6, anchor=center, inner sep=0}, "{d_2}"{description}, from=3-9, to=3-11]
        \arrow[""{name=7, anchor=center, inner sep=0}, "{d_3}"{description}, from=3-11, to=3-13]
        \arrow["{f_{11}}"{description}, between={0.2}{0.8}, squiggly, from=0, to=3]
        \arrow["{f_{12}}"{description}, between={0.2}{0.8}, squiggly, from=0, to=4]
        \arrow["{f_{11}}"{description}, between={0.2}{0.8}, squiggly, from=1, to=5]
        \arrow["{f_{12}}"{description}, between={0.2}{0.8}, squiggly, from=1, to=6]
        \arrow["{f_{33}}"{description}, between={0.2}{0.8}, squiggly, from=2, to=7]
    \end{tikzcd}
\end{equation}


The example on the left corresponds to the tuple $([1] \xra{\langle 02 \rangle} {[2]},\{f_{11},f_{12}\})$ representing a morphism $[1](c_1) \longra [2](d_1,d_2)$ and the one on the right to the tuple $([3] \xra{\langle 0 2 2 3 \rangle} [3],\{f_{11},f_{12},f_{33}\})$ representing a morphism $[3](c_1,c_2,c_3) \longra [3](d_1,d_2,d_3)$.
\end{observation}

\begin{definition}\label{ch3:: def: theta_n}
    Let $n\geq 0$. We define $\bTheta_{1} = \bDelta$ the terminal category, and let 
    \begin{equation}
        \bTheta_{n} = \bDelta \wr \bTheta_{n-1}.
    \end{equation} 
    The objects of $\Theta_{n}$ are given by $[n](c_1,\ldots,c_n)$ where $[n]\in \bDelta$ and $(c_1,\ldots,c_n)\in \bTheta_{n-1}^{\times n}$. 
\end{definition}
Unpacking the definition of $\bTheta_2$ we observe that its objects admit a familiar interpretation as pasting diagrams. As examples, we showcase the following list:
\begin{equation}
\begin{tikzcd}[ampersand replacement=\&]
	{[0]} \& {=} \& \bullet \\
	\\
	{[1]([0])} \& {=} \& \bullet \&\& \bullet \\
	\\
	{[1]([1])} \& {=} \& \bullet \&\& \bullet \\
	\\
	{[1]([2])} \& {=} \& \bullet \&\& \bullet \\
	\\
	{[3]([1],[0],[1])} \& {=} \& \bullet \&\& \bullet \&\& \bullet \&\& \bullet
	\arrow["0"', shift right=5, draw=none, from=1-3, to=1-3, loop, in=55, out=125, distance=10mm]
	\arrow["0"', shift right=5, draw=none, from=3-3, to=3-3, loop, in=55, out=125, distance=10mm]
	\arrow[from=3-3, to=3-5]
	\arrow["1"', shift right=5, draw=none, from=3-5, to=3-5, loop, in=55, out=125, distance=10mm]
	\arrow["0"', shift right=5, draw=none, from=5-3, to=5-3, loop, in=55, out=125, distance=10mm]
	\arrow[""{name=0, anchor=center, inner sep=0}, curve={height=18pt}, from=5-3, to=5-5]
	\arrow[""{name=1, anchor=center, inner sep=0}, curve={height=-18pt}, from=5-3, to=5-5]
	\arrow["1"', shift right=5, draw=none, from=5-5, to=5-5, loop, in=55, out=125, distance=10mm]
	\arrow["0"', shift right=5, draw=none, from=7-3, to=7-3, loop, in=55, out=125, distance=10mm]
	\arrow[""{name=2, anchor=center, inner sep=0}, from=7-3, to=7-5]
	\arrow[""{name=3, anchor=center, inner sep=0}, curve={height=-18pt}, from=7-3, to=7-5]
	\arrow[""{name=4, anchor=center, inner sep=0}, curve={height=18pt}, from=7-3, to=7-5]
	\arrow["1"', shift right=5, draw=none, from=7-5, to=7-5, loop, in=55, out=125, distance=10mm]
	\arrow["0"', shift right=5, draw=none, from=9-3, to=9-3, loop, in=55, out=125, distance=10mm]
	\arrow[""{name=5, anchor=center, inner sep=0}, curve={height=18pt}, from=9-3, to=9-5]
	\arrow[""{name=6, anchor=center, inner sep=0}, curve={height=-18pt}, from=9-3, to=9-5]
	\arrow["1"', shift right=5, draw=none, from=9-5, to=9-5, loop, in=55, out=125, distance=10mm]
	\arrow[from=9-5, to=9-7]
	\arrow["2"', shift right=5, draw=none, from=9-7, to=9-7, loop, in=55, out=125, distance=10mm]
	\arrow[""{name=7, anchor=center, inner sep=0}, curve={height=18pt}, from=9-7, to=9-9]
	\arrow[""{name=8, anchor=center, inner sep=0}, curve={height=-18pt}, from=9-7, to=9-9]
	\arrow["3"', shift right=5, draw=none, from=9-9, to=9-9, loop, in=55, out=125, distance=10mm]
	\arrow[between={0.2}{0.8}, Rightarrow, from=1, to=0]
	\arrow[between={0.2}{0.8}, Rightarrow, from=3, to=2]
	\arrow[between={0.2}{0.8}, Rightarrow, from=2, to=4]
	\arrow[between={0.2}{0.8}, Rightarrow, from=6, to=5]
	\arrow[between={0.2}{0.8}, Rightarrow, from=8, to=7]
\end{tikzcd}
\end{equation}

\medskip 

\nid It is also instructive to make sense of what morphisms in this category do. Since a lot of the intuition for how $\bTheta_n$ works can be drawn from $n=2$, let us focus on this particular case: the object $[0]$ is terminal, so every other object naturally admits a unique morphism to it. Note that these diagrams come endowed with an implicit choice of framing — namely, the 2-cells are directed top to bottom and, of course, the 1-cells are directed left to right. This corresponds to a choice, and does not reflect any inherent geometric properties of the category.

\begin{observation}
    Note that as we iterate the wreath product, its functoriality provides a generalization of the Segal functor for every $n$, given as
    \begin{equation}
        \cT_n = \alpha \circ (\cT \wr \cT_{n-1})\maps \bDelta \wr \bTheta_{n-1} \longra \Gamma \wr \Gamma \longra \Gamma.
    \end{equation}
    This functor, also called an \emph{assembly functor}, sends a pasting diagram to its set of $n$-cells.
\end{observation}

There is an important factorization system on $\bTheta_n$.

\begin{definition}
    Let $\alpha\colon [n] \longra [m]$ in $\bDelta$ be \textbf{inert} if it is the inclusion of a subinterval in $[m]$. More concretely, for every $i$
    \begin{equation}
        \alpha(i+1) = \alpha(0) + i.
    \end{equation}
    Let $(\alpha,{f}_{ij})$ in $\bTheta_n$ be \textbf{inert} if $\alpha$ is inert in $\bDelta$ and every $f_{ij}$ is inert in $\bTheta_{n-1}$. 
\end{definition}

\begin{definition}
    Let $\alpha \colon [n] \longra [m]$ in $\bDelta$ be \textbf{active} if it preserves the endpoints $\alpha(0) = 0$ and $\alpha(n) = m$. As we did for the inert morphisms above, we define active morphisms in $\bTheta_n$ inductively. Let $(\alpha, f_{ij})$ be \textbf{active} if $\alpha$ is active in $\bDelta$ and every $f_{ij}$ is active in $\bTheta_{n-1}$.
\end{definition}


\subsubsection{Segal \texorpdfstring{$\bTheta_n$}{Theta_n}-spaces}

The Segal conditions for $\bTheta_n$-objects say that, similarly to the $n=1$ case, a functor $\cF\maps \bTheta_n \longra \sset$ is Segal if $\cF(T)$ can be computed as the limit of $\cF$ on the basic $i$-cells.

\begin{definition}
    Let us denote by $c_i$, for $i=0,\ldots,n$, the basic $i$-cells in $\bTheta_n$. We define these objects inductively: for $n=1$, let $c_1 = [1]$ and $c_0 = [0]$; for general $n$, let $c_i = [1](c_{i-1})$.
\end{definition}


\begin{definition}
    Let $\GG_n$ called the \textbf{$n$-globular category} denote the full subcategory of $\bTheta_{n,\mathsf{inr}}$ which contains the $i$-cells $c_0,\ldots,c_n$. This subcategory might be depicted as 
    \begin{equation}
        \begin{tikzcd}[ampersand replacement=\&]
            {c_0} \&\& {c_1} \& {} \& \cdots \&\& {c_n}
            \arrow[shift right=2, from=1-1, to=1-3]
            \arrow[shift left=2, from=1-1, to=1-3]
            \arrow[shift right=2, from=1-3, to=1-5]
            \arrow[shift left=2, from=1-3, to=1-5]
            \arrow[shift right=2, from=1-5, to=1-7]
            \arrow[shift left=2, from=1-5, to=1-7]
        \end{tikzcd}
    \end{equation}  
    Let $i_n \maps \GG_n \hookra \bTheta_{n,\mathsf{inr}}$ denote the inclusion of the full $n$-globular subcategory. Given $T\in \bTheta_n$, we define
    \begin{equation}
        \GG_n /T = \GG_n \underset{\bTheta_{n,\mathsf{inr}}}{\times} \bTheta_{n,\mathsf{inr}}/ T.
    \end{equation} 
    The objects of this category correspond bijectively to the set of $k$-cells of $T$ for every $0\leq k \leq n$. 
\end{definition}

\begin{definition}
    A presheaf $\cF\maps \bTheta_{n}^{\sf op} \longra \sset$ is called a \textbf{Segal $\bTheta_n$-object} if the restriction of $\cF$ to its wide subcategory of \textit{inert} morphisms agrees with $\rkan_{i_{n}}(\cF_{|\GG_n})$. More explicitly, the structure map
    \begin{equation}
        \cF(T) \longra \hoLimit[(c_{i} \xra{\mathsf{inrt}} T)\in \GG_{n}^{\sf op}/T] \cF(c_{i})
    \end{equation}
    is an equivalence. 
\end{definition}

Explicitly, this means that if we evaluate $\cF$ on the following diagram 
\begin{equation}
\begin{tikzcd}[ampersand replacement=\&]
	\bullet \&\& \bullet \&\& \bullet
	\arrow["0"', shift right=5, draw=none, from=1-1, to=1-1, loop, in=55, out=125, distance=10mm]
	\arrow[from=1-1, to=1-3]
	\arrow["1"', shift right=5, draw=none, from=1-3, to=1-3, loop, in=55, out=125, distance=10mm]
	\arrow[""{name=0, anchor=center, inner sep=0}, curve={height=-18pt}, from=1-3, to=1-5]
	\arrow[""{name=1, anchor=center, inner sep=0}, curve={height=18pt}, from=1-3, to=1-5]
	\arrow[""{name=2, anchor=center, inner sep=0}, from=1-3, to=1-5]
	\arrow["2"', shift right=5, draw=none, from=1-5, to=1-5, loop, in=55, out=125, distance=10mm]
	\arrow[between={0.2}{0.8}, Rightarrow, from=0, to=2]
	\arrow[between={0.2}{0.8}, Rightarrow, from=2, to=1]
\end{tikzcd}    
\end{equation} 
we can first evaluate it on 
\begin{equation}
    \begin{tikzcd}[ampersand replacement=\&]
        \bullet \&\& \bullet
        \arrow["0"', shift right=5, draw=none, from=1-1, to=1-1, loop, in=55, out=125, distance=10mm]
        \arrow[from=1-1, to=1-3]
        \arrow["1"', shift right=5, draw=none, from=1-3, to=1-3, loop, in=55, out=125, distance=10mm]
    \end{tikzcd} \qquad \text{ and } \qquad
    \begin{tikzcd}[ampersand replacement=\&]
	\bullet \&\& \bullet
	\arrow["1"', shift right=5, draw=none, from=1-1, to=1-1, loop, in=55, out=125, distance=10mm]
	\arrow[""{name=0, anchor=center, inner sep=0}, curve={height=-18pt}, from=1-1, to=1-3]
	\arrow[""{name=1, anchor=center, inner sep=0}, curve={height=18pt}, from=1-1, to=1-3]
	\arrow[""{name=2, anchor=center, inner sep=0}, from=1-1, to=1-3]
	\arrow["2"', shift right=5, draw=none, from=1-3, to=1-3, loop, in=55, out=125, distance=10mm]
	\arrow[between={0.2}{0.8}, Rightarrow, from=0, to=2]
	\arrow[between={0.2}{0.8}, Rightarrow, from=2, to=1]
    \end{tikzcd} 
\end{equation} separately and taking the fiber product. The rightmost piece of the diagram can now be computed by  evaluating 
\begin{equation}
    \begin{tikzcd}[ampersand replacement=\&]
	\bullet \&\& \bullet
	\arrow["1"', shift right=5, draw=none, from=1-1, to=1-1, loop, in=55, out=125, distance=10mm]
	\arrow[""{name=0, anchor=center, inner sep=0}, from=1-1, to=1-3]
	\arrow[""{name=1, anchor=center, inner sep=0}, curve={height=-18pt}, from=1-1, to=1-3]
	\arrow["2"', shift right=5, draw=none, from=1-3, to=1-3, loop, in=55, out=125, distance=10mm]
	\arrow[between={0.2}{0.8}, Rightarrow, from=1, to=0]
\end{tikzcd} \qquad \text{ and } \qquad
\begin{tikzcd}[ampersand replacement=\&]
	\bullet \&\& \bullet
	\arrow["1"', shift right=5, draw=none, from=1-1, to=1-1, loop, in=55, out=125, distance=10mm]
	\arrow[""{name=0, anchor=center, inner sep=0}, from=1-1, to=1-3]
	\arrow[""{name=1, anchor=center, inner sep=0}, curve={height=18pt}, from=1-1, to=1-3]
	\arrow["2"', shift right=5, draw=none, from=1-3, to=1-3, loop, in=55, out=125, distance=10mm]
	\arrow[between={0.2}{0.8}, Rightarrow, from=0, to=1]
\end{tikzcd}
\end{equation} separately and, again, further taking the fiber product.


\subsection{Category of higher cubes}
Before delving into the definition of $\cube_n$, the main character of this section, let us recall that the interval category $\Int$ admits a canonical functor 
\begin{equation}
    \gamma_1 \colon \Int^{\op} \longra \Gamma
\end{equation}
which is defined as the composition of $\gamma^{\op} \colon \Int^{\op} \longra \bDelta$ and the Segal functor $\cT \colon \bDelta \longra \Gamma$. It sends an interval $I$ to the set of path components of its complement $[0,1]\setminus I$. Predicated on the existence of this functor, we introduce the notion of the $n$-cubes category, a higher dimensional adaptation of $\Int$. 

The definition below, including the iterated wreath-product presentation and the intended localization to $\bTheta_n$, is based on unpublished work of Manuel Araújo and Pedro Boavida de Brito.

\begin{definition}\label{def:: ch6: cube_n category wreath}
    Let $\cube_1^{\op} = \Int^{\op}$. We define the \define{$n$-cubes category} to be the iterated wreath product of $\Int$ with itself $n$ times: 
    \begin{equation}
        \cube_n^{\op} = \cube_1^{\op} \wr \cube_{n-1}^{\op}. 
    \end{equation}
\end{definition}

Alternatively, we can describe the objects of $\cube_n$ as follows: let $[0,1]^{n}$ denote the $n$-dimensional unit cube. The objects of $\cube_n$ can be described as open subsets $R \subseteq [0,1]^{n}$ consisting of:
    \begin{enumerate}
        \item An open collar \[C \simeq [0,\epsilon_l)\times [0,1]^{n-1} \coprod (\epsilon_r,1]\times [0,1]^{n-1},\] of $\partial [0,1] \times [0,1]^{n-1}$, where $0<\epsilon_l< \epsilon_r < 1$.
        \item In the complement of the collar, $R$ consists of a disjoint union of open cubes of the following form:
        \begin{equation}
            \coprod_{\alpha \in A} I_1^{\alpha}\times J_{\alpha}
        \end{equation}
        where $A$ is the set of inner connected components of $I_1 \in \Int_1$ and such that 
        \begin{equation}
            C\simeq I_1^{\sf out}\times [0,1]^{n-1} .
        \end{equation}  
        and $J_{\alpha} \in \cube_{n-1}$.
    \end{enumerate}
Morphisms are given by inclusions of open subsets.


\begin{remark}\label{ch6:: rmk: unpack cube_n}
    We now unpack the definition. For $n=0$, we have that $\cube_0 = \ast$ by definition. For $n=1$, objects are merely finite disjoint unions of open intervals of $[0,1]$ which always must include a collar of the boundary $\{0,1\}\subset [0,1]$. In other words, $\cube_1 = \Int$, i.e., the interval category. For $n=2$, things get immediately more interesting. As a way of illustrating the geometric conditions forced by condition (2), we depict below three examples of finite disjoint unions of open cubes in $[0,1]^2$: 
    \begin{equation}
        \tikzfig{figures/FO/REC1}{\begin{tikzpicture}
	\begin{pgfonlayer}{nodelayer}
		\node [style=none] (109) at (-11.025, -3.5) {};
		\node [style=none] (110) at (-11.025, -4.5) {};
		\node [style=none] (111) at (-9.025, -4.5) {};
		\node [style=none] (112) at (-9.025, -3.5) {};
		\node [style=none] (113) at (-11, 4.5) {};
		\node [style=none] (114) at (-9, 4.5) {};
		\node [style=none] (115) at (-11, 3.5) {};
		\node [style=none] (116) at (-9, 3.5) {};
		\node [style=none] (105) at (-13, -3.5) {};
		\node [style=none] (106) at (-13, -4.5) {};
		\node [style=none] (107) at (-12, -4.5) {};
		\node [style=none] (108) at (-12, -3.5) {};
		\node [style=none] (101) at (-13, 4.5) {};
		\node [style=none] (102) at (-13, 3.5) {};
		\node [style=none] (103) at (-12, 3.5) {};
		\node [style=none] (104) at (-12, 4.5) {};
		\node [style=none] (0) at (-16, 4.5) {};
		\node [style=none] (1) at (-16, -4.5) {};
		\node [style=none] (26) at (-14.5, 4.5) {};
		\node [style=none] (27) at (-14.5, -4.5) {};
		\node [style=none] (2) at (-7, -4.5) {};
		\node [style=none] (3) at (-7, 4.5) {};
		\node [style=none] (4) at (-13, 2) {};
		\node [style=none] (5) at (-13, 1) {};
		\node [style=none] (6) at (-12, 1) {};
		\node [style=none] (7) at (-12, 2) {};
		\node [style=none] (8) at (-13, -1) {};
		\node [style=none] (9) at (-13, -3) {};
		\node [style=none] (10) at (-12, -3) {};
		\node [style=none] (11) at (-12, -1) {};
		\node [style=none] (12) at (-11, 3) {};
		\node [style=none] (13) at (-11, 2) {};
		\node [style=none] (14) at (-9, 3) {};
		\node [style=none] (15) at (-9, 2) {};
		\node [style=none] (16) at (-11, 1) {};
		\node [style=none] (17) at (-11, -2) {};
		\node [style=none] (18) at (-9, 1) {};
		\node [style=none] (19) at (-9, -2) {};
		\node [style=none] (20) at (-9, -2) {};
		\node [style=none] (21) at (-11, -2.5) {};
		\node [style=none] (22) at (-11, -2.75) {};
		\node [style=none] (23) at (-9, -2.5) {};
		\node [style=none] (24) at (-9, -2.75) {};
		\node [style=none] (25) at (-9, -2.75) {};
		\node [style=none] (28) at (-8, 4.5) {};
		\node [style=none] (29) at (-8, -4.5) {};
		\node [style=none] (30) at (-7, 4.5) {};
		\node [style=none] (31) at (-7, -4.5) {};
		\node [style=none] (32) at (-17, -5.5) {};
		\node [style=none] (33) at (-16.25, -5.5) {};
		\node [style=none] (34) at (-17, -4.75) {};
		\node [style=none] (35) at (-16.25, -6) {$1$};
		\node [style=none] (36) at (-17.5, -4.75) {$2$};
		\node [style=none] (37) at (-4.5, 4.5) {};
		\node [style=none] (38) at (-4.5, -4.5) {};
		\node [style=none] (39) at (-3, 4.5) {};
		\node [style=none] (40) at (-3, -4.5) {};
		\node [style=none] (41) at (4.5, -4.5) {};
		\node [style=none] (42) at (4.5, 4.5) {};
		\node [style=none] (43) at (-1.25, 2.75) {};
		\node [style=none] (44) at (-1.25, 1.75) {};
		\node [style=none] (45) at (0.25, 1.75) {};
		\node [style=none] (46) at (0.25, 2.75) {};
		\node [style=none] (47) at (-1.5, -1) {};
		\node [style=none] (48) at (-1.5, -3) {};
		\node [style=none] (49) at (-0.5, -3) {};
		\node [style=none] (50) at (-0.5, -1) {};
		\node [style=none] (51) at (0.5, 3) {};
		\node [style=none] (52) at (0.5, 2) {};
		\node [style=none] (53) at (2.5, 3) {};
		\node [style=none] (54) at (2.5, 2) {};
		\node [style=none] (55) at (0.25, 1) {};
		\node [style=none] (56) at (0.25, -2) {};
		\node [style=none] (57) at (1.25, 1) {};
		\node [style=none] (58) at (1.25, -2) {};
		\node [style=none] (59) at (1.25, -2) {};
		\node [style=none] (60) at (0.5, -2.5) {};
		\node [style=none] (61) at (0.5, -2.75) {};
		\node [style=none] (62) at (2.5, -2.5) {};
		\node [style=none] (63) at (2.5, -2.75) {};
		\node [style=none] (64) at (2.5, -2.75) {};
		\node [style=none] (65) at (3.5, 4.5) {};
		\node [style=none] (66) at (3.5, -4.5) {};
		\node [style=none] (67) at (4.5, 4.5) {};
		\node [style=none] (68) at (4.5, -4.5) {};
		\node [style=none] (69) at (7, 4.5) {};
		\node [style=none] (70) at (7, -4.5) {};
		\node [style=none] (71) at (8.5, 4.5) {};
		\node [style=none] (72) at (8.5, -4.5) {};
		\node [style=none] (73) at (16, -4.5) {};
		\node [style=none] (74) at (16, 4.5) {};
		\node [style=none] (87) at (10.25, 1) {};
		\node [style=none] (88) at (10.25, -2) {};
		\node [style=none] (89) at (12.25, 1) {};
		\node [style=none] (90) at (12.25, -2) {};
		\node [style=none] (91) at (12.25, -2) {};
		\node [style=none] (97) at (13.5, 4.5) {};
		\node [style=none] (98) at (13.5, -4.5) {};
		\node [style=none] (99) at (16, 4.5) {};
		\node [style=none] (100) at (16, -4.5) {};
		\node [style=none] (117) at (18, -5.5) {};
	\end{pgfonlayer}
	\begin{pgfonlayer}{edgelayer}
		\draw [style=cube_1] (27.center)
			 to (1.center)
			 to (0.center)
			 to (26.center)
			 to cycle;
		\draw [style=geometric] (1.center) to (2.center);
		\draw [style=geometric] (2.center) to (3.center);
		\draw [style=geometric, in=360, out=180] (3.center) to (0.center);
		\draw [style=cube_1] (7.center)
			 to (4.center)
			 to (5.center)
			 to (6.center)
			 to cycle;
		\draw [style=cube_1] (8.center)
			 to (9.center)
			 to (10.center)
			 to (11.center)
			 to cycle;
		\draw [style=cube_1] (12.center)
			 to (13.center)
			 to (15.center)
			 to (14.center)
			 to cycle;
		\draw [style=cube_1] (18.center)
			 to (16.center)
			 to (17.center)
			 to (20.center)
			 to cycle;
		\draw [style=cube_1] (23.center)
			 to (21.center)
			 to (22.center)
			 to (25.center)
			 to cycle;
		\draw [style=cube_1] (31.center)
			 to (29.center)
			 to (28.center)
			 to (30.center)
			 to cycle;
		\draw [style=geometric] (0.center) to (1.center);
		\draw [style=cd-ar] (32.center) to (33.center);
		\draw [style=cd-ar] (32.center) to (34.center);
		\draw [style=cube_1] (40.center)
			 to (38.center)
			 to (37.center)
			 to (39.center)
			 to cycle;
		\draw [style=geometric] (38.center) to (41.center);
		\draw [style=cube_1] (66.center)
			 to (65.center)
			 to (67.center)
			 to (68.center)
			 to cycle;
		\draw [style=geometric, in=360, out=180] (42.center) to (37.center);
		\draw [style=cube_1] (46.center)
			 to (43.center)
			 to (44.center)
			 to (45.center)
			 to cycle;
		\draw [style=cube_1] (47.center)
			 to (48.center)
			 to (49.center)
			 to (50.center)
			 to cycle;
		\draw [style=cube_1] (51.center)
			 to (52.center)
			 to (54.center)
			 to (53.center)
			 to cycle;
		\draw [style=cube_1] (57.center)
			 to (55.center)
			 to (56.center)
			 to (59.center)
			 to cycle;
		\draw [style=cube_1] (61.center)
			 to (64.center)
			 to (62.center)
			 to (60.center)
			 to cycle;
		\draw [style=geometric] (37.center) to (38.center);
		\draw [style=cube_1] (72.center)
			 to (70.center)
			 to (69.center)
			 to (71.center)
			 to cycle;
		\draw [style=geometric] (70.center) to (73.center);
		\draw [style=cube_1] (98.center)
			 to (97.center)
			 to (99.center)
			 to (100.center)
			 to cycle;
		\draw [style=geometric, in=360, out=180] (74.center) to (69.center);
		\draw [style=cube_1] (89.center)
			 to (87.center)
			 to (88.center)
			 to (91.center)
			 to cycle;
		\draw [style=geometric] (69.center) to (70.center);
		\draw [style=geometric] (67.center) to (68.center);
		\draw [style=geometric] (100.center) to (99.center);
		\draw [style=cube_1] (104.center)
			 to (101.center)
			 to (102.center)
			 to (103.center)
			 to cycle;
		\draw [style=cube_1] (108.center)
			 to (105.center)
			 to (106.center)
			 to (107.center)
			 to cycle;
		\draw [style=cube_1] (112.center)
			 to (109.center)
			 to (110.center)
			 to (111.center)
			 to cycle;
		\draw [style=cube_1] (115.center)
			 to (116.center)
			 to (114.center)
			 to (113.center)
			 to cycle;
	\end{pgfonlayer}
\end{tikzpicture}
}
    \end{equation}
    only the left one actually corresponds to an object of $\cube_2$, the one on the right fails to satisfy the collar condition and the one in the middle fails to satisfy both the collar and the product-decomposition condition. To indicate a morphism we might make use of the opacity as a way of distinguishing objects. The next figure illustrates this in a convenient manner, by drawing the overlapping subsets on top of the morphism arrow:
    \begin{equation}
        \tikzfig{figures/FO/REC4}{\begin{tikzpicture}
	\begin{pgfonlayer}{nodelayer}
		\node [style=none] (109) at (0.475, -3.5) {};
		\node [style=none] (110) at (0.475, -4.5) {};
		\node [style=none] (111) at (2.475, -4.5) {};
		\node [style=none] (112) at (2.475, -3.5) {};
		\node [style=none] (113) at (0.5, 4.5) {};
		\node [style=none] (114) at (2.5, 4.5) {};
		\node [style=none] (115) at (0.5, 3.5) {};
		\node [style=none] (116) at (2.5, 3.5) {};
		\node [style=none] (105) at (-1.5, -3.5) {};
		\node [style=none] (106) at (-1.5, -4.5) {};
		\node [style=none] (107) at (-0.5, -4.5) {};
		\node [style=none] (108) at (-0.5, -3.5) {};
		\node [style=none] (101) at (-1.5, 4.5) {};
		\node [style=none] (102) at (-1.5, 3.5) {};
		\node [style=none] (103) at (-0.5, 3.5) {};
		\node [style=none] (104) at (-0.5, 4.5) {};
		\node [style=none] (0) at (-4.5, 4.5) {};
		\node [style=none] (1) at (-4.5, -4.5) {};
		\node [style=none] (26) at (-3, 4.5) {};
		\node [style=none] (27) at (-3, -4.5) {};
		\node [style=none] (2) at (4.5, -4.5) {};
		\node [style=none] (3) at (4.5, 4.5) {};
		\node [style=none] (4) at (-1.5, 2) {};
		\node [style=none] (5) at (-1.5, 1) {};
		\node [style=none] (6) at (-0.5, 1) {};
		\node [style=none] (7) at (-0.5, 2) {};
		\node [style=none] (8) at (-1.5, -0.5) {};
		\node [style=none] (9) at (-1.5, -2.5) {};
		\node [style=none] (10) at (-0.5, -2.5) {};
		\node [style=none] (11) at (-0.5, -0.5) {};
		\node [style=none] (12) at (0.5, 3) {};
		\node [style=none] (13) at (0.5, 2) {};
		\node [style=none] (14) at (2.5, 3) {};
		\node [style=none] (15) at (2.5, 2) {};
		\node [style=none] (16) at (0.5, 1) {};
		\node [style=none] (17) at (0.5, -2) {};
		\node [style=none] (18) at (2.5, 1) {};
		\node [style=none] (19) at (2.5, -2) {};
		\node [style=none] (20) at (2.5, -2) {};
		\node [style=none] (21) at (0.5, -2.5) {};
		\node [style=none] (22) at (0.5, -2.75) {};
		\node [style=none] (23) at (2.5, -2.5) {};
		\node [style=none] (24) at (2.5, -2.75) {};
		\node [style=none] (25) at (2.5, -2.75) {};
		\node [style=none] (28) at (3.5, 4.5) {};
		\node [style=none] (29) at (3.5, -4.5) {};
		\node [style=none] (30) at (4.5, 4.5) {};
		\node [style=none] (31) at (4.5, -4.5) {};
		\node [style=none] (32) at (-5.5, -5.5) {};
		\node [style=none] (33) at (-4.75, -5.5) {};
		\node [style=none] (34) at (-5.5, -4.75) {};
		\node [style=none] (35) at (-4.75, -6) {$1$};
		\node [style=none] (36) at (-6, -4.75) {$2$};
		\node [style=none] (117) at (-1.75, 2.25) {};
	\end{pgfonlayer}
	\begin{pgfonlayer}{edgelayer}
		\draw [style={cube_1}] (27.center)
			 to (1.center)
			 to (0.center)
			 to (26.center)
			 to cycle;
		\draw [style=geometric] (1.center) to (2.center);
		\draw [style=geometric] (2.center) to (3.center);
		\draw [style=geometric, in=360, out=180] (3.center) to (0.center);
		\draw [style={cube_1}] (7.center)
			 to (4.center)
			 to (5.center)
			 to (6.center)
			 to cycle;
		\draw [style={cube_1}] (8.center)
			 to (9.center)
			 to (10.center)
			 to (11.center)
			 to cycle;
		\draw [style={cube_1}] (12.center)
			 to (13.center)
			 to (15.center)
			 to (14.center)
			 to cycle;
		\draw [style={cube_1}] (18.center)
			 to (16.center)
			 to (17.center)
			 to (20.center)
			 to cycle;
		\draw [style={cube_1}] (23.center)
			 to (21.center)
			 to (22.center)
			 to (25.center)
			 to cycle;
		\draw [style={cube_1}] (31.center)
			 to (29.center)
			 to (28.center)
			 to (30.center)
			 to cycle;
		\draw [style=geometric] (0.center) to (1.center);
		\draw [style=cd-ar] (32.center) to (33.center);
		\draw [style=cd-ar] (32.center) to (34.center);
		\draw [style={cube_1}] (104.center)
			 to (101.center)
			 to (102.center)
			 to (103.center)
			 to cycle;
		\draw [style={cube_1}] (108.center)
			 to (105.center)
			 to (106.center)
			 to (107.center)
			 to cycle;
		\draw [style={cube_1}] (112.center)
			 to (109.center)
			 to (110.center)
			 to (111.center)
			 to cycle;
		\draw [style={cube_1}] (115.center)
			 to (116.center)
			 to (114.center)
			 to (113.center)
			 to cycle;
	\end{pgfonlayer}
\end{tikzpicture}
} \,\,\,\,\,\, \xra{\scalebox{0.6}{\tikzfig{figures/FO/REC2}{\begin{tikzpicture}
	\begin{pgfonlayer}{nodelayer}
		\node [style=none] (109) at (0.475, -3.5) {};
		\node [style=none] (110) at (0.475, -4.5) {};
		\node [style=none] (111) at (2.475, -4.5) {};
		\node [style=none] (112) at (2.475, -3.5) {};
		\node [style=none] (113) at (0.5, 4.5) {};
		\node [style=none] (114) at (2.5, 4.5) {};
		\node [style=none] (115) at (0.5, 3.5) {};
		\node [style=none] (116) at (2.5, 3.5) {};
		\node [style=none] (105) at (-1.5, -3.5) {};
		\node [style=none] (106) at (-1.5, -4.5) {};
		\node [style=none] (107) at (-0.5, -4.5) {};
		\node [style=none] (108) at (-0.5, -3.5) {};
		\node [style=none] (101) at (-1.5, 4.5) {};
		\node [style=none] (102) at (-1.5, 3.5) {};
		\node [style=none] (103) at (-0.5, 3.5) {};
		\node [style=none] (104) at (-0.5, 4.5) {};
		\node [style=none] (0) at (-4.5, 4.5) {};
		\node [style=none] (1) at (-4.5, -4.5) {};
		\node [style=none] (26) at (-3, 4.5) {};
		\node [style=none] (27) at (-3, -4.5) {};
		\node [style=none] (2) at (4.5, -4.5) {};
		\node [style=none] (3) at (4.5, 4.5) {};
		\node [style=none] (4) at (-1.5, 2) {};
		\node [style=none] (5) at (-1.5, 1) {};
		\node [style=none] (6) at (-0.5, 1) {};
		\node [style=none] (7) at (-0.5, 2) {};
		\node [style=none] (8) at (-1.5, -0.5) {};
		\node [style=none] (9) at (-1.5, -2.5) {};
		\node [style=none] (10) at (-0.5, -2.5) {};
		\node [style=none] (11) at (-0.5, -0.5) {};
		\node [style=none] (12) at (0.5, 3) {};
		\node [style=none] (13) at (0.5, 2) {};
		\node [style=none] (14) at (2.5, 3) {};
		\node [style=none] (15) at (2.5, 2) {};
		\node [style=none] (16) at (0.5, 1) {};
		\node [style=none] (17) at (0.5, -2) {};
		\node [style=none] (18) at (2.5, 1) {};
		\node [style=none] (19) at (2.5, -2) {};
		\node [style=none] (20) at (2.5, -2) {};
		\node [style=none] (21) at (0.5, -2.5) {};
		\node [style=none] (22) at (0.5, -2.75) {};
		\node [style=none] (23) at (2.5, -2.5) {};
		\node [style=none] (24) at (2.5, -2.75) {};
		\node [style=none] (25) at (2.5, -2.75) {};
		\node [style=none] (28) at (3.5, 4.5) {};
		\node [style=none] (29) at (3.5, -4.5) {};
		\node [style=none] (30) at (4.5, 4.5) {};
		\node [style=none] (31) at (4.5, -4.5) {};
		\node [style=none] (117) at (-1.75, 2.25) {};
		\node [style=none] (119) at (-2, -3) {};
		\node [style=none] (120) at (0, -3) {};
		\node [style=none] (121) at (-2, 2.5) {};
		\node [style=none] (122) at (0, 2.5) {};
		\node [style=none] (124) at (-2, -3.5) {};
		\node [style=none] (125) at (-2, -4.5) {};
		\node [style=none] (126) at (0, -3.5) {};
		\node [style=none] (127) at (0, -4.5) {};
		\node [style=none] (128) at (-2, 4.5) {};
		\node [style=none] (129) at (0, 4.5) {};
		\node [style=none] (130) at (-2, 3) {};
		\node [style=none] (131) at (0, 3) {};
		\node [style=none] (132) at (0, -6) {};
	\end{pgfonlayer}
	\begin{pgfonlayer}{edgelayer}
		\draw [style={cube_1}] (27.center)
			 to (1.center)
			 to (0.center)
			 to (26.center)
			 to cycle;
		\draw [style=geometric] (1.center) to (2.center);
		\draw [style=geometric] (2.center) to (3.center);
		\draw [style=geometric, in=360, out=180] (3.center) to (0.center);
		\draw [style={cube_1}] (7.center)
			 to (4.center)
			 to (5.center)
			 to (6.center)
			 to cycle;
		\draw [style={cube_1}] (8.center)
			 to (9.center)
			 to (10.center)
			 to (11.center)
			 to cycle;
		\draw [style={cube_1}] (12.center)
			 to (13.center)
			 to (15.center)
			 to (14.center)
			 to cycle;
		\draw [style={cube_1}] (18.center)
			 to (16.center)
			 to (17.center)
			 to (20.center)
			 to cycle;
		\draw [style={cube_1}] (23.center)
			 to (21.center)
			 to (22.center)
			 to (25.center)
			 to cycle;
		\draw [style={cube_1}] (31.center)
			 to (29.center)
			 to (28.center)
			 to (30.center)
			 to cycle;
		\draw [style=geometric] (0.center) to (1.center);
		\draw [style={cube_1}] (104.center)
			 to (101.center)
			 to (102.center)
			 to (103.center)
			 to cycle;
		\draw [style={cube_1}] (108.center)
			 to (105.center)
			 to (106.center)
			 to (107.center)
			 to cycle;
		\draw [style={cube_1}] (112.center)
			 to (109.center)
			 to (110.center)
			 to (111.center)
			 to cycle;
		\draw [style={cube_1}] (115.center)
			 to (116.center)
			 to (114.center)
			 to (113.center)
			 to cycle;
		\draw [style={cube_2}] (120.center)
			 to (122.center)
			 to (121.center)
			 to (119.center)
			 to cycle;
		\draw [style={cube_2}] (126.center)
			 to (124.center)
			 to (125.center)
			 to (127.center)
			 to cycle;
		\draw [style={cube_2}] (129.center)
			 to (131.center)
			 to (130.center)
			 to (128.center);
	\end{pgfonlayer}
\end{tikzpicture}
}}} \,\,\,\,\,\,\, \tikzfig{figures/FO/REC3}{\begin{tikzpicture}
	\begin{pgfonlayer}{nodelayer}
		\node [style=none] (109) at (0.475, -3.5) {};
		\node [style=none] (110) at (0.475, -4.5) {};
		\node [style=none] (111) at (2.475, -4.5) {};
		\node [style=none] (112) at (2.475, -3.5) {};
		\node [style=none] (113) at (0.5, 4.5) {};
		\node [style=none] (114) at (2.5, 4.5) {};
		\node [style=none] (115) at (0.5, 3.5) {};
		\node [style=none] (116) at (2.5, 3.5) {};
		\node [style=none] (0) at (-4.5, 4.5) {};
		\node [style=none] (1) at (-4.5, -4.5) {};
		\node [style=none] (26) at (-3, 4.5) {};
		\node [style=none] (27) at (-3, -4.5) {};
		\node [style=none] (2) at (4.5, -4.5) {};
		\node [style=none] (3) at (4.5, 4.5) {};
		\node [style=none] (12) at (0.5, 3) {};
		\node [style=none] (13) at (0.5, 2) {};
		\node [style=none] (14) at (2.5, 3) {};
		\node [style=none] (15) at (2.5, 2) {};
		\node [style=none] (16) at (0.5, 1) {};
		\node [style=none] (17) at (0.5, -2) {};
		\node [style=none] (18) at (2.5, 1) {};
		\node [style=none] (19) at (2.5, -2) {};
		\node [style=none] (20) at (2.5, -2) {};
		\node [style=none] (21) at (0.5, -2.5) {};
		\node [style=none] (22) at (0.5, -2.75) {};
		\node [style=none] (23) at (2.5, -2.5) {};
		\node [style=none] (24) at (2.5, -2.75) {};
		\node [style=none] (25) at (2.5, -2.75) {};
		\node [style=none] (28) at (3.5, 4.5) {};
		\node [style=none] (29) at (3.5, -4.5) {};
		\node [style=none] (30) at (4.5, 4.5) {};
		\node [style=none] (31) at (4.5, -4.5) {};
		\node [style=none] (32) at (5.75, -5.5) {};
		\node [style=none] (33) at (6.5, -5.5) {};
		\node [style=none] (34) at (5.75, -4.75) {};
		\node [style=none] (35) at (6.5, -6) {$1$};
		\node [style=none] (36) at (5.25, -4.75) {$2$};
		\node [style=none] (117) at (-1.75, 2.25) {};
		\node [style=none] (119) at (-2, -3) {};
		\node [style=none] (120) at (0, -3) {};
		\node [style=none] (121) at (-2, 2.5) {};
		\node [style=none] (122) at (0, 2.5) {};
		\node [style=none] (124) at (-2, -3.5) {};
		\node [style=none] (125) at (-2, -4.5) {};
		\node [style=none] (126) at (0, -3.5) {};
		\node [style=none] (127) at (0, -4.5) {};
		\node [style=none] (128) at (-2, 4.5) {};
		\node [style=none] (129) at (0, 4.5) {};
		\node [style=none] (130) at (-2, 3) {};
		\node [style=none] (131) at (0, 3) {};
	\end{pgfonlayer}
	\begin{pgfonlayer}{edgelayer}
		\draw [style={cube_1}] (27.center)
			 to (1.center)
			 to (0.center)
			 to (26.center)
			 to cycle;
		\draw [style=geometric] (1.center) to (2.center);
		\draw [style=geometric] (2.center) to (3.center);
		\draw [style=geometric, in=360, out=180] (3.center) to (0.center);
		\draw [style={cube_1}] (12.center)
			 to (13.center)
			 to (15.center)
			 to (14.center)
			 to cycle;
		\draw [style={cube_1}] (18.center)
			 to (16.center)
			 to (17.center)
			 to (20.center)
			 to cycle;
		\draw [style={cube_1}] (23.center)
			 to (21.center)
			 to (22.center)
			 to (25.center)
			 to cycle;
		\draw [style={cube_1}] (31.center)
			 to (29.center)
			 to (28.center)
			 to (30.center)
			 to cycle;
		\draw [style=geometric] (0.center) to (1.center);
		\draw [style=cd-ar] (32.center) to (33.center);
		\draw [style=cd-ar] (32.center) to (34.center);
		\draw [style={cube_1}] (112.center)
			 to (109.center)
			 to (110.center)
			 to (111.center)
			 to cycle;
		\draw [style={cube_1}] (115.center)
			 to (116.center)
			 to (114.center)
			 to (113.center)
			 to cycle;
		\draw [style={cube_1}] (120.center)
			 to (122.center)
			 to (121.center)
			 to (119.center)
			 to cycle;
		\draw [style={cube_1}] (126.center)
			 to (124.center)
			 to (125.center)
			 to (127.center)
			 to cycle;
		\draw [style={cube_1}] (129.center)
			 to (131.center)
			 to (130.center)
			 to (128.center);
	\end{pgfonlayer}
\end{tikzpicture}
}
    \end{equation}
\end{remark}


\noindent We now initiate the comparison between $\cube_n$ and $\bTheta_n$, analogously to the way we defined $\gamma\colon \Int \longra \bDelta^{\op}$ in Chapter 2. In fact, our definition will be just as iterative as our definition of $\cube_n$.

One of the advantages of defining $\cube_n$ using an iterated wreath product construction, is that it automatically comes with a functor to $\bTheta_{n}^{\op}$ — defined inductively as follows:
    \begin{equation}
        \gamma_n = \gamma_{1}\wr \gamma_{n-1} \maps \cube_n \longra \bTheta_n^{\op}.
    \end{equation}
This maps 
    \begin{equation}
        \coprod_{\alpha \in A} I_1^{\alpha}\times J_{\alpha} \,\,\longmapsto \,\,\gamma_1(I_1)\big(\gamma_{n-1}(J_1),\ldots, \gamma_{n-1}(J_k)\big).
    \end{equation}

\begin{remark}
    The functor $\gamma_n$, as is defined, might seem opaque and devoid of immediate geometric content, yet the underlying idea underpinning it is wholly geometric. To sharpen our discussion and in order to undress the definition of its inductive nature, let us think about the case of $\gamma_2$. Let us consider the object $R$ in $\cube_2$ which we depicted above:
    \begin{equation}
    \tikzfig{figures/FO/REC4}{\begin{tikzpicture}
	\begin{pgfonlayer}{nodelayer}
		\node [style=none] (109) at (0.475, -3.5) {};
		\node [style=none] (110) at (0.475, -4.5) {};
		\node [style=none] (111) at (2.475, -4.5) {};
		\node [style=none] (112) at (2.475, -3.5) {};
		\node [style=none] (113) at (0.5, 4.5) {};
		\node [style=none] (114) at (2.5, 4.5) {};
		\node [style=none] (115) at (0.5, 3.5) {};
		\node [style=none] (116) at (2.5, 3.5) {};
		\node [style=none] (105) at (-1.5, -3.5) {};
		\node [style=none] (106) at (-1.5, -4.5) {};
		\node [style=none] (107) at (-0.5, -4.5) {};
		\node [style=none] (108) at (-0.5, -3.5) {};
		\node [style=none] (101) at (-1.5, 4.5) {};
		\node [style=none] (102) at (-1.5, 3.5) {};
		\node [style=none] (103) at (-0.5, 3.5) {};
		\node [style=none] (104) at (-0.5, 4.5) {};
		\node [style=none] (0) at (-4.5, 4.5) {};
		\node [style=none] (1) at (-4.5, -4.5) {};
		\node [style=none] (26) at (-3, 4.5) {};
		\node [style=none] (27) at (-3, -4.5) {};
		\node [style=none] (2) at (4.5, -4.5) {};
		\node [style=none] (3) at (4.5, 4.5) {};
		\node [style=none] (4) at (-1.5, 2) {};
		\node [style=none] (5) at (-1.5, 1) {};
		\node [style=none] (6) at (-0.5, 1) {};
		\node [style=none] (7) at (-0.5, 2) {};
		\node [style=none] (8) at (-1.5, -0.5) {};
		\node [style=none] (9) at (-1.5, -2.5) {};
		\node [style=none] (10) at (-0.5, -2.5) {};
		\node [style=none] (11) at (-0.5, -0.5) {};
		\node [style=none] (12) at (0.5, 3) {};
		\node [style=none] (13) at (0.5, 2) {};
		\node [style=none] (14) at (2.5, 3) {};
		\node [style=none] (15) at (2.5, 2) {};
		\node [style=none] (16) at (0.5, 1) {};
		\node [style=none] (17) at (0.5, -2) {};
		\node [style=none] (18) at (2.5, 1) {};
		\node [style=none] (19) at (2.5, -2) {};
		\node [style=none] (20) at (2.5, -2) {};
		\node [style=none] (21) at (0.5, -2.5) {};
		\node [style=none] (22) at (0.5, -2.75) {};
		\node [style=none] (23) at (2.5, -2.5) {};
		\node [style=none] (24) at (2.5, -2.75) {};
		\node [style=none] (25) at (2.5, -2.75) {};
		\node [style=none] (28) at (3.5, 4.5) {};
		\node [style=none] (29) at (3.5, -4.5) {};
		\node [style=none] (30) at (4.5, 4.5) {};
		\node [style=none] (31) at (4.5, -4.5) {};
		\node [style=none] (32) at (-5.5, -5.5) {};
		\node [style=none] (33) at (-4.75, -5.5) {};
		\node [style=none] (34) at (-5.5, -4.75) {};
		\node [style=none] (35) at (-4.75, -6) {$1$};
		\node [style=none] (36) at (-6, -4.75) {$2$};
		\node [style=none] (117) at (-1.75, 2.25) {};
	\end{pgfonlayer}
	\begin{pgfonlayer}{edgelayer}
		\draw [style={cube_1}] (27.center)
			 to (1.center)
			 to (0.center)
			 to (26.center)
			 to cycle;
		\draw [style=geometric] (1.center) to (2.center);
		\draw [style=geometric] (2.center) to (3.center);
		\draw [style=geometric, in=360, out=180] (3.center) to (0.center);
		\draw [style={cube_1}] (7.center)
			 to (4.center)
			 to (5.center)
			 to (6.center)
			 to cycle;
		\draw [style={cube_1}] (8.center)
			 to (9.center)
			 to (10.center)
			 to (11.center)
			 to cycle;
		\draw [style={cube_1}] (12.center)
			 to (13.center)
			 to (15.center)
			 to (14.center)
			 to cycle;
		\draw [style={cube_1}] (18.center)
			 to (16.center)
			 to (17.center)
			 to (20.center)
			 to cycle;
		\draw [style={cube_1}] (23.center)
			 to (21.center)
			 to (22.center)
			 to (25.center)
			 to cycle;
		\draw [style={cube_1}] (31.center)
			 to (29.center)
			 to (28.center)
			 to (30.center)
			 to cycle;
		\draw [style=geometric] (0.center) to (1.center);
		\draw [style=cd-ar] (32.center) to (33.center);
		\draw [style=cd-ar] (32.center) to (34.center);
		\draw [style={cube_1}] (104.center)
			 to (101.center)
			 to (102.center)
			 to (103.center)
			 to cycle;
		\draw [style={cube_1}] (108.center)
			 to (105.center)
			 to (106.center)
			 to (107.center)
			 to cycle;
		\draw [style={cube_1}] (112.center)
			 to (109.center)
			 to (110.center)
			 to (111.center)
			 to cycle;
		\draw [style={cube_1}] (115.center)
			 to (116.center)
			 to (114.center)
			 to (113.center)
			 to cycle;
	\end{pgfonlayer}
\end{tikzpicture}
} 
    \end{equation}
    This object can be written as 
    \begin{equation}
       R = (I_1^{1} \times J_1)\, \coprod\, (I_{1}^{2} \times J_2)
    \end{equation}
    such that
    \begin{align}
        I_1 & = \tikzfig{figures/FO/REC5}{\begin{tikzpicture}
	\begin{pgfonlayer}{nodelayer}
		\node [style=none] (32) at (-5, -1) {};
		\node [style=none] (33) at (-4.25, -1) {};
		\node [style=none] (35) at (-4.25, -1.5) {$1$};
		\node [style=none] (118) at (-0.025, 0) {};
		\node [style=none] (119) at (1.975, 0) {};
		\node [style=none] (120) at (-2, 0) {};
		\node [style=none] (121) at (-1, 0) {};
		\node [style=none] (122) at (-5, 0) {};
		\node [style=none] (123) at (-3.5, 0) {};
		\node [style=none] (124) at (4, 0) {};
		\node [style=none] (125) at (3, 0) {};
		\node [style=none] (126) at (4, 0) {};
		\node [style=endpoint] (127) at (4, 0) {};
		\node [style=endpoint] (128) at (-5, 0) {};
	\end{pgfonlayer}
	\begin{pgfonlayer}{edgelayer}
		\draw [style=cd-ar] (32.center) to (33.center);
		\draw [style={cube_1}] (123.center) to (122.center);
		\draw [style=none] (122.center) to (124.center);
		\draw [style={cube_1}] (126.center) to (125.center);
		\draw [style={cube_1}] (120.center) to (121.center);
		\draw [style={cube_1}] (118.center) to (119.center);
		\draw [style={open_int}] (122.center) to (123.center);
		\draw [style={open_int}] (120.center) to (121.center);
		\draw [style={open_int}] (118.center) to (119.center);
		\draw [style={open_int}] (125.center) to (126.center);
	\end{pgfonlayer}
\end{tikzpicture}
}\,, \\ 
        J_1 & = \tikzfig{figures/FO/REC6}{\begin{tikzpicture}
	\begin{pgfonlayer}{nodelayer}
		\node [style=none] (32) at (-5, -1) {};
		\node [style=none] (33) at (-4.25, -1) {};
		\node [style=none] (35) at (-4.25, -1.5) {$2$};
		\node [style=none] (122) at (-5, 0) {};
		\node [style=none] (124) at (4, 0) {};
		\node [style=none] (126) at (4, 0) {};
		\node [style=endpoint] (127) at (4, 0) {};
		\node [style=endpoint] (128) at (-5, 0) {};
		\node [style=none] (129) at (-4, 0) {};
		\node [style=none] (130) at (-5, 0) {};
		\node [style=none] (131) at (4, 0) {};
		\node [style=none] (132) at (3, 0) {};
		\node [style=none] (133) at (1.5, 0) {};
		\node [style=none] (134) at (0.5, 0) {};
		\node [style=none] (135) at (-1, 0) {};
		\node [style=none] (136) at (-3, 0) {};
	\end{pgfonlayer}
	\begin{pgfonlayer}{edgelayer}
		\draw [style=cd-ar] (32.center) to (33.center);
		\draw [style=none] (122.center) to (124.center);
		\draw [style={open_int}] (133.center) to (134.center);
		\draw [style={open_int}] (135.center) to (136.center);
		\draw [style={open_int}] (131.center) to (132.center);
		\draw [style={open_int}] (129.center) to (130.center);
	\end{pgfonlayer}
\end{tikzpicture}
}\,, \\ 
        J_2 &= \tikzfig{figures/FO/REC7}{\begin{tikzpicture}
	\begin{pgfonlayer}{nodelayer}
		\node [style=none] (32) at (-5, -1) {};
		\node [style=none] (33) at (-4.25, -1) {};
		\node [style=none] (35) at (-4.25, -1.5) {$2$};
		\node [style=none] (122) at (-5, 0) {};
		\node [style=none] (124) at (4, 0) {};
		\node [style=none] (126) at (4, 0) {};
		\node [style=endpoint] (127) at (4, 0) {};
		\node [style=endpoint] (128) at (-5, 0) {};
		\node [style=none] (130) at (-5, 0) {};
		\node [style=none] (131) at (4, 0) {};
		\node [style=none] (132) at (-4, 0) {};
		\node [style=none] (133) at (-5, 0) {};
		\node [style=none] (134) at (4, 0) {};
		\node [style=none] (135) at (3, 0) {};
		\node [style=none] (136) at (2.5, 0) {};
		\node [style=none] (137) at (1.5, 0) {};
		\node [style=none] (138) at (0.5, 0) {};
		\node [style=none] (139) at (-2.5, 0) {};
		\node [style=none] (140) at (-3, 0) {};
		\node [style=none] (141) at (-3.25, 0) {};
	\end{pgfonlayer}
	\begin{pgfonlayer}{edgelayer}
		\draw [style=cd-ar] (32.center) to (33.center);
		\draw [style={cube_1}] (138.center) to (139.center);
		\draw [style={cube_1}] (140.center) to (141.center);
		\draw [style={cube_1}] (132.center) to (133.center);
		\draw [style=none] (122.center) to (124.center);
		\draw [style={cube_1}] (136.center) to (137.center);
		\draw [style={cube_1}] (134.center) to (135.center);
		\draw [style={open_int}] (133.center) to (132.center);
		\draw [style={open_int}] (140.center) to (141.center);
		\draw [style={open_int}] (139.center) to (138.center);
		\draw [style={open_int}] (137.center) to (136.center);
		\draw [style={open_int}] (135.center) to (134.center);
	\end{pgfonlayer}
\end{tikzpicture}
}\,. 
    \end{align}
    with the arrows at the bottom highlighting the directions along which $\gamma_2$ assigns $[2]([2],[3])$ to $R$, where the left outer $[2]$ computes $\pi_0\big([0,1]\setminus\pr_1(R)\big)$ with the linear order inherited from the standard orientation on $[0,1]$. The left inner $[2]$ corresponds to $\gamma_{1}(J_1)$ and the right inner $[3]$ corresponds to $\gamma_1(J_2)$. Interpreting this object as a pasting diagram, we note that it corresponds to the following one:
    \begin{equation}
    \begin{tikzcd}
        \bullet && \bullet && \bullet
        \arrow[""{name=0, anchor=center, inner sep=0}, from=1-1, to=1-3]
        \arrow[""{name=1, anchor=center, inner sep=0}, curve={height=-18pt}, from=1-1, to=1-3]
        \arrow[""{name=2, anchor=center, inner sep=0}, curve={height=18pt}, from=1-1, to=1-3]
        \arrow[""{name=3, anchor=center, inner sep=0}, curve={height=-30pt}, from=1-3, to=1-5]
        \arrow[""{name=4, anchor=center, inner sep=0}, curve={height=30pt}, from=1-3, to=1-5]
        \arrow[""{name=5, anchor=center, inner sep=0}, curve={height=12pt}, from=1-3, to=1-5]
        \arrow[""{name=6, anchor=center, inner sep=0}, curve={height=-12pt}, from=1-3, to=1-5]
        \arrow[between={0.2}{0.8}, Rightarrow, from=1, to=0]
        \arrow[between={0.2}{0.8}, Rightarrow, from=0, to=2]
        \arrow[between={0.2}{0.8}, Rightarrow, from=3, to=6]
        \arrow[between={0.2}{0.8}, Rightarrow, from=6, to=5]
        \arrow[between={0.2}{0.8}, Rightarrow, from=5, to=4]
    \end{tikzcd}
    \end{equation}
    This interpretation of the objects of $\cube_n$ is an important one. 

    To give this example some additional flesh, we study also the action of $\gamma_2$ under the morphism depicted in \Cref{ch6:: rmk: unpack cube_n}. The morphism is induced by an inclusion along the first direction which induces a $\pi_0$-isomorphism corresponding to a slight thickening of the first interval component, and two inclusions along the second direction: the first merges the two cubes over the first component whilst the second is just an equality. As seen from the image of $\gamma_2$, this corresponds to the morphism 
    \begin{equation}
        \begin{tikzcd}
            \bullet && \bullet && \bullet
            \arrow[""{name=0, anchor=center, inner sep=0}, curve={height=-18pt}, from=1-1, to=1-3]
            \arrow[""{name=1, anchor=center, inner sep=0}, curve={height=18pt}, from=1-1, to=1-3]
            \arrow[""{name=2, anchor=center, inner sep=0}, curve={height=-30pt}, from=1-3, to=1-5]
            \arrow[""{name=3, anchor=center, inner sep=0}, curve={height=30pt}, from=1-3, to=1-5]
            \arrow[""{name=4, anchor=center, inner sep=0}, curve={height=12pt}, from=1-3, to=1-5]
            \arrow[""{name=5, anchor=center, inner sep=0}, curve={height=-12pt}, from=1-3, to=1-5]
            \arrow[between={0.2}{0.8}, Rightarrow, from=0, to=1]
            \arrow[between={0.2}{0.8}, Rightarrow, from=2, to=5]
            \arrow[between={0.2}{0.8}, Rightarrow, from=5, to=4]
            \arrow[between={0.2}{0.8}, Rightarrow, from=4, to=3]
        \end{tikzcd} \longrightarrow 
        \begin{tikzcd}
            \bullet && \bullet && \bullet
            \arrow[""{name=0, anchor=center, inner sep=0}, from=1-1, to=1-3]
            \arrow[""{name=1, anchor=center, inner sep=0}, curve={height=-18pt}, from=1-1, to=1-3]
            \arrow[""{name=2, anchor=center, inner sep=0}, curve={height=18pt}, from=1-1, to=1-3]
            \arrow[""{name=3, anchor=center, inner sep=0}, curve={height=-30pt}, from=1-3, to=1-5]
            \arrow[""{name=4, anchor=center, inner sep=0}, curve={height=30pt}, from=1-3, to=1-5]
            \arrow[""{name=5, anchor=center, inner sep=0}, curve={height=12pt}, from=1-3, to=1-5]
            \arrow[""{name=6, anchor=center, inner sep=0}, curve={height=-12pt}, from=1-3, to=1-5]
            \arrow[between={0.2}{0.8}, Rightarrow, from=1, to=0]
            \arrow[between={0.2}{0.8}, Rightarrow, from=0, to=2]
            \arrow[between={0.2}{0.8}, Rightarrow, from=3, to=6]
            \arrow[between={0.2}{0.8}, Rightarrow, from=6, to=5]
            \arrow[between={0.2}{0.8}, Rightarrow, from=5, to=4]
        \end{tikzcd}
    \end{equation}  
    which vertically composes the two cells on the left, also encoded as
    \begin{equation}
        [2]([1],[3]) \xra{\id (\langle 0,2 \rangle, \id)} [2]([2],[3])
    \end{equation}
\end{remark}


\begin{remark}
    Although we will mostly use the geometric model, it is helpful to keep in mind the inductive form of morphisms. A morphism
    \begin{equation}
        f\maps I_{1}(J_1,\ldots, J_n) \longra I'_{1}(J'_{1},\ldots,J'_{k})
    \end{equation}
    in $\cube_n$ may be viewed as an inclusion in the first coordinate together with, for each inner connected component of $I_1$, a morphism in $\cube_{n-1}$ from the corresponding $J_i$ to the lower-dimensional piece of the target lying over the image component. In the geometric model this is simply the inclusion of one rectangular decomposition into another.
\end{remark}

\subsection{Segal \texorpdfstring{$n$}{n}-cube-spaces}
The guiding idea is the same as in Chapter 2. Rather than work with all morphisms in $\Int = \cube_1$, we singled out those which ought to become invertible once the model is topologized, and localized along them. For $\cube_n$ we follow the same philosophy: we isolate the morphisms which should play the role of weak equivalences from the geometric point of view.

\begin{remark}
For $n=1$, these are precisely the inclusions $I \longra J$ of open intervals inducing isomorphisms on $\pi_0$. 
\end{remark}

We generalize this set as follows.

\begin{definition}
    Let $\cW_{\pi_{0}}(n)$ denote the subcategory of $\cube_n$ spanned by the morphisms 
    \begin{equation}
    f\colon R \longra S, \quad \text{ such that } \gamma_n(f) = \id.
    \end{equation} 
\end{definition}

This is the higher-dimensional analogue of the class $\cW_{\pi_{0}}$ introduced in Chapter 2. Geometrically, one should think of the morphisms in $\cW_{\pi_{0}}(n)$ as those inclusions obtained by enlarging the cubes of $R$ without creating new overlaps or forcing previously distinct rectangles to merge. This suggests the following conjectural comparison with the usual $\bTheta_n$-model.   

\begin{conjecture} \label{fo:: conj: local.}
    There should be a Quillen equivalence
\begin{equation}
    \begin{tikzcd}[ampersand replacement=\&]
	{\fun(\cube_{n},\sset)_{\mathsf{inj}}} \&\& {\fun(\bTheta_{n}^{\op},\sset)_{\mathsf{inj}}.}
	\arrow[""{name=0, anchor=center, inner sep=0}, "{\gamma_{n}^*}", bend left=10, from=1-1, to=1-3]
	\arrow[""{name=1, anchor=center, inner sep=0}, "{\rkan_{\gamma_{n}}}", bend left=10, from=1-3, to=1-1]
	\arrow["\dashv"{anchor=center, rotate=-90}, draw=none, from=1, to=0]
\end{tikzcd}
\end{equation}
induced by the functor $\gamma_{n} \maps \cube_{n} \longra \bTheta_{n}^{\op}$.
\end{conjecture}

We only indicate why one might expect this to hold.

\begin{proof}[Sketch of the expected argument]
    The blueprint of this proof is identical to the one for $n=1$ (see the proof of \Cref{ch2:: prop: weak contractibility of pi_0}), with an important caveat: that proof relies on the contractibility of ordered configuration spaces on the line. For higher $n$, this ceases to be true even for $n=2$. Evidently, this means we need to adapt our previous argument. We use instead a subspace $\conf^{\sf special}_T \subset ([0,1]^n)^{\pi_{0}(R_{T})}$ consisting of configurations adapted to the tree-like structure of $T$. This is inspired by Ayala--Hepworth's description of configuration-space cells indexed by $n$-orderings, or equivalently by healthy planar level trees \cite{ayalaConfigurationSpaces$Theta_n$2014}; in the present rectangular setting, one expects $\conf^{\sf special}_T$ to be a convex cell.
    \begin{align}
        |\cJ_T| \,\, &\simeq \,\, \hoColimit[R_T \in \cJ_T]\ast \\
        &\overset{(1)}{\simeq} \,\,\hoColimit[R_T \in \cJ_T] \conf(R_T) \cap \conf^{\sf special}_T \\
        &\overset{(2)}{\simeq} \,\,\conf^{\sf special}_T \\
        &\overset{(3)}{\simeq} \,\,\ast 
    \end{align}

    Some remarks about how this proof works: both $\conf(R_T)$ and $\conf^{\sf special}_T$ are contractible spaces; the former being homeomorphic to a product of Euclidean spaces and the latter, by design, to a convex cell in $([0,1]^{n})^{\pi_0(R_T)}$, thereby rendering their intersection contractible, too; this proves $(1)$. To obtain $(2)$, one argues as in the interval case: given a special configuration $\phi \in \conf^{\sf special}_T$, the subposet 
    \begin{equation}
        \big\{R_T \in \cJ_T\mkern1mu :\, \phi \in \conf(R_T)\cap \conf^{\sf special}_T\big\} \subset \cJ_{T}
    \end{equation}
    is filtered, therefore it has contractible geometric realization. Lurie's $\oo$-categorical version of Seifert-Van Kampen then identifies the homotopy colimit of the cover by the contractible pieces $\conf(R_T)\cap \conf^{\sf special}_T$ with $\conf^{\sf special}_T$. Finally, $(3)$ follows from the convexity of $\conf^{\sf special}_T$.
\end{proof}

Assuming \Cref{fo:: conj: local.}, we will call a functor $\cF \maps \cube_n \longra \sset$ a $\cube_n$-space if it sends morphisms in $\cW_{\pi_{0}}(n)$ to weak equivalences. We will call it a \emph{Segal} $\cube_n$-\emph{space} if, in addition, for every $R \in \cube_n$, the value of $\cF(R)$ can be recovered by successively cutting $R$ into single cubes, that is, into the pieces representing $k$-cells, one coordinate direction at a time and in increasing order. A typical such decomposition is illustrated by
\begin{equation}
    \tikzfig{figures/FO/REC14}{\begin{tikzpicture}
	\begin{pgfonlayer}{nodelayer}
		\node [style=none] (109) at (0.225, -3.5) {};
		\node [style=none] (110) at (0.225, -4.5) {};
		\node [style=none] (111) at (2.225, -4.5) {};
		\node [style=none] (112) at (2.225, -3.5) {};
		\node [style=none] (113) at (0.25, 4.5) {};
		\node [style=none] (114) at (2.25, 4.5) {};
		\node [style=none] (115) at (0.25, 3.5) {};
		\node [style=none] (116) at (2.25, 3.5) {};
		\node [style=none] (107) at (-0.75, -4.5) {};
		\node [style=none] (108) at (-0.75, -3.5) {};
		\node [style=none] (103) at (-0.75, 3.5) {};
		\node [style=none] (104) at (-0.75, 4.5) {};
		\node [style=none] (0) at (-4.75, 4.5) {};
		\node [style=none] (1) at (-4.75, -4.5) {};
		\node [style=none] (26) at (-0.75, 4.5) {};
		\node [style=none] (27) at (-0.75, -4.5) {};
		\node [style=none] (2) at (4.25, -4.5) {};
		\node [style=none] (3) at (4.25, 4.5) {};
		\node [style=none] (6) at (-0.75, 1) {};
		\node [style=none] (7) at (-0.75, 2) {};
		\node [style=none] (10) at (-0.75, -2.5) {};
		\node [style=none] (11) at (-0.75, -0.5) {};
		\node [style=none] (16) at (0.25, 1) {};
		\node [style=none] (17) at (0.25, -2) {};
		\node [style=none] (18) at (2.25, 1) {};
		\node [style=none] (19) at (2.25, -2) {};
		\node [style=none] (20) at (2.25, -2) {};
		\node [style=none] (24) at (2.25, -2.75) {};
		\node [style=none] (28) at (3.25, 4.5) {};
		\node [style=none] (29) at (3.25, -4.5) {};
		\node [style=none] (30) at (4.25, 4.5) {};
		\node [style=none] (31) at (4.25, -4.5) {};
		\node [style=none] (117) at (-8.525, 1) {};
		\node [style=none] (118) at (-8.525, 2) {};
		\node [style=none] (119) at (-12.5, 10) {};
		\node [style=none] (120) at (-8.5, 10) {};
		\node [style=none] (121) at (-12.5, 1) {};
		\node [style=none] (122) at (-8.5, 1) {};
		\node [style=none] (123) at (-14.5, 2) {};
		\node [style=none] (124) at (-14.5, 1) {};
		\node [style=none] (125) at (-13.5, 1) {};
		\node [style=none] (126) at (-13.5, 2) {};
		\node [style=none] (127) at (-14.5, 10) {};
		\node [style=none] (128) at (-13.5, 10) {};
		\node [style=none] (129) at (-17.5, 10) {};
		\node [style=none] (130) at (-17.5, 1) {};
		\node [style=none] (131) at (-16, 10) {};
		\node [style=none] (132) at (-16, 1) {};
		\node [style=none] (133) at (-8.5, 1) {};
		\node [style=none] (134) at (-8.5, 10) {};
		\node [style=none] (135) at (-14.5, 10) {};
		\node [style=none] (136) at (-14.5, 6.5) {};
		\node [style=none] (137) at (-13.5, 6.5) {};
		\node [style=none] (138) at (-13.5, 10) {};
		\node [style=none] (139) at (-14.5, 5) {};
		\node [style=none] (140) at (-14.5, 3) {};
		\node [style=none] (141) at (-13.5, 3) {};
		\node [style=none] (142) at (-13.5, 5) {};
		\node [style=none] (143) at (-8.5, 6.5) {};
		\node [style=none] (144) at (-8.5, 3.5) {};
		\node [style=none] (145) at (-8.5, 3.5) {};
		\node [style=none] (146) at (-8.5, 2.75) {};
		\node [style=none] (147) at (-8.5, 10) {};
		\node [style=none] (148) at (-8.5, 1) {};
		\node [style=none] (149) at (-14.75, 10.25) {};
		\node [style=none] (150) at (-8.525, -10) {};
		\node [style=none] (151) at (-8.525, -9) {};
		\node [style=none] (152) at (-12.5, -1) {};
		\node [style=none] (153) at (-8.5, -1) {};
		\node [style=none] (154) at (-12.5, -10) {};
		\node [style=none] (155) at (-8.5, -10) {};
		\node [style=none] (156) at (-14.5, -9) {};
		\node [style=none] (157) at (-14.5, -10) {};
		\node [style=none] (158) at (-13.5, -10) {};
		\node [style=none] (159) at (-13.5, -9) {};
		\node [style=none] (160) at (-14.5, -1) {};
		\node [style=none] (161) at (-13.5, -1) {};
		\node [style=none] (162) at (-17.5, -1) {};
		\node [style=none] (163) at (-17.5, -10) {};
		\node [style=none] (164) at (-16, -1) {};
		\node [style=none] (165) at (-16, -10) {};
		\node [style=none] (166) at (-8.5, -10) {};
		\node [style=none] (167) at (-8.5, -1) {};
		\node [style=none] (168) at (-14.5, -1) {};
		\node [style=none] (169) at (-14.5, -4.5) {};
		\node [style=none] (170) at (-13.5, -4.5) {};
		\node [style=none] (171) at (-13.5, -1) {};
		\node [style=none] (172) at (-14.5, -6) {};
		\node [style=none] (173) at (-14.5, -8) {};
		\node [style=none] (174) at (-13.5, -8) {};
		\node [style=none] (175) at (-13.5, -6) {};
		\node [style=none] (176) at (-8.5, -4.5) {};
		\node [style=none] (177) at (-8.5, -7.5) {};
		\node [style=none] (178) at (-8.5, -7.5) {};
		\node [style=none] (179) at (-8.5, -8.25) {};
		\node [style=none] (180) at (-8.5, -1) {};
		\node [style=none] (181) at (-8.5, -10) {};
		\node [style=none] (182) at (-14.75, -0.75) {};
		\node [style=none] (183) at (-13, 0) {$\circ_{2}$};
		\node [style=none] (184) at (-18, 10.75) {};
		\node [style=none] (185) at (-18.75, 10.75) {};
		\node [style=none] (186) at (-18.75, -10.75) {};
		\node [style=none] (187) at (-18, -10.75) {};
		\node [style=none] (188) at (-8.25, 10.75) {};
		\node [style=none] (189) at (-7.5, 10.75) {};
		\node [style=none] (190) at (-7.5, -10.75) {};
		\node [style=none] (191) at (-8.25, -10.75) {};
		\node [style=none] (192) at (-6.25, 0) {$\circ_{1}$};
	\end{pgfonlayer}
	\begin{pgfonlayer}{edgelayer}
		\draw [style={cube_1}] (27.center)
			 to (1.center)
			 to (0.center)
			 to (26.center)
			 to cycle;
		\draw [style=geometric] (1.center) to (2.center);
		\draw [style=geometric] (2.center) to (3.center);
		\draw [style=geometric, in=360, out=180] (3.center) to (0.center);
		\draw [style={cube_1}] (6.center) to (7.center);
		\draw [style={cube_1}] (10.center) to (11.center);
		\draw [style={cube_1}] (18.center)
			 to (16.center)
			 to (17.center)
			 to (20.center)
			 to cycle;
		\draw [style={cube_1}] (31.center)
			 to (29.center)
			 to (28.center)
			 to (30.center)
			 to cycle;
		\draw [style=geometric] (0.center) to (1.center);
		\draw [style={cube_1}] (103.center) to (104.center);
		\draw [style={cube_1}] (107.center) to (108.center);
		\draw [style={cube_1}] (112.center)
			 to (109.center)
			 to (110.center)
			 to (111.center)
			 to cycle;
		\draw [style={cube_1}] (115.center)
			 to (116.center)
			 to (114.center)
			 to (113.center)
			 to cycle;
		\draw [style={cube_1}] (132.center)
			 to (130.center)
			 to (129.center)
			 to (131.center)
			 to cycle;
		\draw [style=geometric] (130.center) to (133.center);
		\draw [style=geometric] (133.center) to (134.center);
		\draw [style=geometric, in=360, out=180] (134.center) to (129.center);
		\draw [style={cube_1}] (138.center)
			 to (135.center)
			 to (136.center)
			 to (137.center)
			 to cycle;
		\draw [style={cube_1}] (139.center)
			 to (140.center)
			 to (141.center)
			 to (142.center)
			 to cycle;
		\draw [style={cube_1}] (145.center) to (143.center);
		\draw [style={cube_1}] (147.center) to (148.center);
		\draw [style=geometric] (129.center) to (130.center);
		\draw [style={cube_1}] (128.center) to (127.center);
		\draw [style={cube_1}] (126.center)
			 to (123.center)
			 to (124.center)
			 to (125.center)
			 to cycle;
		\draw [style={cube_1}] (117.center) to (118.center);
		\draw [style={cube_1}] (121.center)
			 to (122.center)
			 to (120.center)
			 to (119.center)
			 to cycle;
		\draw [style={cube_1}] (165.center)
			 to (163.center)
			 to (162.center)
			 to (164.center)
			 to cycle;
		\draw [style=geometric] (163.center) to (166.center);
		\draw [style=geometric] (166.center) to (167.center);
		\draw [style=geometric, in=360, out=180] (167.center) to (162.center);
		\draw [style={cube_1}] (171.center)
			 to (168.center)
			 to (169.center)
			 to (170.center)
			 to cycle;
		\draw [style={cube_1}] (172.center)
			 to (173.center)
			 to (174.center)
			 to (175.center)
			 to cycle;
		\draw [style={cube_1}] (178.center) to (176.center);
		\draw [style={cube_1}] (180.center) to (181.center);
		\draw [style=geometric] (162.center) to (163.center);
		\draw [style={cube_1}] (161.center) to (160.center);
		\draw [style={cube_1}] (159.center)
			 to (156.center)
			 to (157.center)
			 to (158.center)
			 to cycle;
		\draw [style={cube_1}] (150.center) to (151.center);
		\draw [style={cube_1}] (154.center)
			 to (155.center)
			 to (153.center)
			 to (152.center)
			 to cycle;
		\draw [style=geometric] (184.center) to (185.center);
		\draw [style=geometric] (185.center) to (186.center);
		\draw [style=geometric] (186.center) to (187.center);
		\draw [style=geometric] (188.center) to (189.center);
		\draw [style=geometric] (189.center) to (190.center);
		\draw [style=geometric] (190.center) to (191.center);
	\end{pgfonlayer}
\end{tikzpicture}
}
\end{equation}
\medskip

\nid Since $\cF$ sends axis-aligned enlargements to equivalences, many different cube configurations may encode the same $k$-cell of $\bTheta_n$. Herein lies the advantage of working with $\cube_n$ rather than directly with $\bTheta_n^{\op}$ in geometric contexts: one keeps the rigid globular bookkeeping, while retaining the flexibility of rectangular decompositions and the freedom to isotope configurations.   

\subsubsection{\texorpdfstring{$(\oo,d)$}{(infty,d)}-categories of stratified cobordisms}

One of the goals of defining $(\oo,d)$-categories using Segal $\cube_d$-spaces is so that we can construct $(\oo,d)$-categories of fully extended stratified bordisms by slightly adapting the methods used in Chapter 3 and Chapter 5. 



\medskip

\nid Before proceeding with this construction, let us first illustrate why naively copying the structure of $\tilbordstr_d$ will fail if we do not begin by addressing the definition of $\cube_n$. The issue is that a general object of $\cube_2$ allows different components in the first direction to carry unrelated decompositions in the second direction. After cutting once, there is then no canonical lower-dimensional cube recording where the next cuts should occur, so the inductive description of composition loses control of the slicing data. Assume $(\oo,d=2)$, and that we are working with 2-dimensional 2-bordisms.  


\begin{equation}
    \tikzfig{figures/Stcob/SCO11}{\begin{tikzpicture}
	\begin{pgfonlayer}{nodelayer}
		\node [style=none] (107) at (-7, 2.25) {};
		\node [style=none] (108) at (-4.5, 5) {};
		\node [style=none] (109) at (-5.25, 2.75) {};
		\node [style=none] (110) at (0.5, 5) {};
		\node [style=none] (111) at (-2, 2.25) {};
		\node [style=none] (112) at (-4.5, 2.75) {};
		\node [style=none] (113) at (-3.5, 4) {};
		\node [style=endpoint] (40) at (-5, -1) {};
		\node [style=none] (48) at (-5, -1) {};
		\node [style=none] (4) at (-5, -1) {};
		\node [style=none] (49) at (-3.25, -1) {};
		\node [style=none] (50) at (-10, -6) {};
		\node [style=none] (51) at (-8.25, -6) {};
		\node [style=none] (5) at (3.75, -1) {};
		\node [style=none] (6) at (-3.25, -1) {};
		\node [style=none] (7) at (2, -1) {};
		\node [style=none] (8) at (-10, -6) {};
		\node [style=none] (9) at (-1.25, -6) {};
		\node [style=none] (10) at (-8.25, -6) {};
		\node [style=none] (11) at (-3, -6) {};
		\node [style=endpoint] (41) at (-10, -6) {};
		\node [style=endpoint] (42) at (-1.25, -6) {};
		\node [style=endpoint] (43) at (3.75, -1) {};
		\node [style=none] (44) at (2, -1) {};
		\node [style=none] (45) at (3.75, -1) {};
		\node [style=none] (52) at (2, -1) {};
		\node [style=none] (53) at (3.75, -1) {};
		\node [style=none] (54) at (-3, -6) {};
		\node [style=none] (55) at (-1.25, -6) {};
		\node [style=none] (56) at (-10, -6) {};
		\node [style=none] (57) at (-5.75, -4.75) {};
		\node [style=none] (58) at (-2.75, -4.75) {};
		\node [style=none] (59) at (-7, -6) {};
		\node [style=none] (60) at (-4, -6) {};
		\node [style=none] (69) at (-2, -1) {};
		\node [style=none] (70) at (0.75, -1) {};
		\node [style=none] (71) at (-3.25, -2.25) {};
		\node [style=none] (72) at (-0.5, -2.25) {};
		\node [style=none] (77) at (-3.75, -2.75) {};
		\node [style=none] (78) at (-1, -2.75) {};
		\node [style=none] (79) at (-5, -4) {};
		\node [style=none] (80) at (-2.25, -4) {};
		\node [style=none] (81) at (-7, -4.75) {};
		\node [style=none] (82) at (-1.75, -4.75) {};
		\node [style=none] (83) at (-4.5, -2.25) {};
		\node [style=none] (84) at (0.75, -2.25) {};
		\node [style=none] (85) at (-7, 5.25) {};
		\node [style=none] (86) at (-1.75, 5.25) {};
		\node [style=none] (87) at (-4.5, -2.25) {};
		\node [style=none] (88) at (0.75, -2.25) {};
		\node [style=none] (89) at (-4.5, 7.75) {};
		\node [style=none] (90) at (0.75, 7.75) {};
		\node [style=none] (92) at (-7, -3.25) {};
		\node [style=none] (95) at (-4.5, -1) {};
		\node [style=none] (100) at (-7, 3.75) {};
		\node [style=none] (101) at (-5.25, 2.75) {};
		\node [style=none] (102) at (-3.5, 4) {};
		\node [style=none] (103) at (-4.5, 2.75) {};
		\node [style=none] (104) at (-1.75, 3.5) {};
		\node [style=none] (105) at (0.5, 6.25) {};
		\node [style=none] (106) at (-4.5, 6.25) {};
		\node [style=none] (114) at (-3.5, -3.25) {};
		\node [style=0 strat] (115) at (-5.725, 3.35) {};
		\node [style=0 strat] (116) at (-5.75, 2.325) {};
		\node [style=0 strat] (118) at (-3.25, 4.75) {};
		\node [style=0 strat] (119) at (-1.75, 6.25) {};
		\node [style=0 strat] (122) at (-4.75, 4.5) {};
		\node [style=none] (123) at (-4.75, 3.375) {};
	\end{pgfonlayer}
	\begin{pgfonlayer}{edgelayer}
		\draw [style=ts-dashed-grey] (89.center) to (87.center);
		\draw [style=geometric shaded] (108.center)
			 to (110.center)
			 to (111.center)
			 to [in=-120, out=-180, looseness=0.75] (112.center)
			 to [in=-15, out=60] (113.center)
			 to [in=45, out=165] (109.center)
			 to [in=0, out=-90, looseness=0.50] (107.center)
			 to cycle;
		\draw [style=box, in=60, out=180] (123.center) to (116);
		\draw [style=ts-dashed-red] (102.center) to (114.center);
		\draw [style=geometric] (4.center) to (5.center);
		\draw [style={cube_1}] (51.center)
			 to (49.center)
			 to (48.center)
			 to (50.center)
			 to cycle;
		\draw [style={open_int}] (4.center) to (6.center);
		\draw [style={open_int}] (7.center) to (5.center);
		\draw [style=geometric] (8.center) to (9.center);
		\draw [style={open_int}] (8.center) to (10.center);
		\draw [style={open_int}] (11.center) to (9.center);
		\draw [style=red shade low op] (44.center) to (45.center);
		\draw [style={cube_1}] (55.center)
			 to (53.center)
			 to (52.center)
			 to (54.center)
			 to cycle;
		\draw [style={cube_1}] (60.center)
			 to (58.center)
			 to (57.center)
			 to (59.center)
			 to cycle;
		\draw [style={cube_1}] (72.center)
			 to (70.center)
			 to (69.center)
			 to (71.center)
			 to cycle;
		\draw [style={cube_1}] (80.center)
			 to (78.center)
			 to (77.center)
			 to (79.center)
			 to cycle;
		\draw [style={open_int}] (59.center) to (60.center);
		\draw [style={open_int}] (69.center) to (70.center);
		\draw [style=ts-dashed-grey] (81.center) to (82.center);
		\draw [style=ts-dashed-grey] (83.center) to (84.center);
		\draw [style=ts-dashed-grey] (84.center) to (82.center);
		\draw [style=ts-dashed-grey] (83.center) to (81.center);
		\draw [style=ts-dashed-grey] (85.center) to (81.center);
		\draw [style=ts-dashed-grey] (85.center) to (86.center);
		\draw [style=ts-dashed-grey] (86.center) to (82.center);
		\draw [style=ts-dashed-grey] (87.center) to (88.center);
		\draw [style=ts-dashed-grey] (89.center) to (90.center);
		\draw [style=ts-dashed-grey] (90.center) to (88.center);
		\draw [style=ts-dashed-grey] (85.center) to (89.center);
		\draw [style=ts-dashed-grey] (86.center) to (90.center);
		\draw [style=geometric shaded] (102.center)
			 to [in=45, out=165] (101.center)
			 to [in=-15, out=90, looseness=0.75] (100.center)
			 to (106.center)
			 to (105.center)
			 to (104.center)
			 to [in=60, out=180, looseness=0.75] (103.center)
			 to [in=-15, out=60] cycle;
		\draw [style=box, in=30, out=-30, looseness=0.75] (122) to (123.center);
		\draw [style=box, in=180, out=45] (122) to (118);
		\draw [style=box] (115) to (122);
		\draw [style=box, in=-150, out=30, looseness=1.25] (118) to (119);
	\end{pgfonlayer}
\end{tikzpicture}
}
\end{equation}
\medskip

The point of the picture is precisely that the two pieces want different interval data in the transverse direction, so they do not naturally refine a common lower-dimensional cutting pattern.

In order to accommodate for the construction of the $(\oo,d)$-category of stratified bordisms, we must restrict $\cube_n$ to a full subcategory $\cR\cube_{n}$ spanned by the objects $X \in \cube_n$ such that

\begin{equation}
    X = I_1 (J_1, \ldots, J_n), \quad \gamma_{1}(I_1) = [n]
\end{equation}
and there exists $J \in \cR\cube_{n-1}$ together with refinements $J_{i} \ra J$ for each $1 \leq i \leq n$. 



\begin{construction}\label{fo:: def: higher strat bord}
     Fix $V$ to be a finite-dimensional real vector space. We would like to define a functor 
    \begin{equation}
    \tilbordstr_{[0,\ldots,d];V}\maps \man^{\op}\times \cR\cube_{d} \longrightarrow \set    
    \end{equation}
    sending a pair $(M,R)$, where $M  \in \man^{\op}$ and $R\in \cR\cube_d$ such that 
    \begin{equation}
        R = I(J_1, \ldots, J_n),
    \end{equation} 
    to the set of bounded closed subspaces 
    \begin{equation}
   \big( W\subseteq M\times \core(R) \times V \big)
    \end{equation} 
    such that the following conditions are satisfied:
    \begin{enumerate}
    \item \label{id: cor 1} $W$ is endowed with the structure of a conically smooth stratification and a $\SMfd_{d}^{\lag d \rag}$-framing. In particular, it is a conically smooth stratified space with a $\lag d \rag$-corner structure
    \begin{equation}
         (W \ra P),\quad (W \xra{c} \cP(\{1,\cdots,d\}));
    \end{equation} 
    and the embedding must preserve the underlying stratifications i.e., send the boundary of $W$ to the boundary of $M\times \core(R) \times V$.
    \item \label{id: cor 2} The restricted projection 
    \begin{equation}
        \pi\maps \big(W \cap (M\times I^{\mathsf{c}}\times [0,1]^{d-1} \times V) \big) \hookra M\times I^{\mathsf{c}}\times [0,1]^{d-1} \times V \fib  M\times I^{\mathsf{c}}\times [0,1]^{d-1}
    \end{equation} 
    is a conically smooth submersion.
    \item \label{id: cor 3} When restricting $W$ to $J_{i} \times I^{\sf inn}$, for every $t\in I^{\sf inn}$ one expects the resulting slice over $t$ to determine an element of $\tilbordstr_{[0,\ldots,d-1];V}(M,J_{i})$.
    \item \label{id: cor 4} The restricted projection 
    \begin{equation}
    \pi_W\maps W \hookra M\times \core(R) \times V \fib M 
    \end{equation}
    is a conically smooth submersion.
    \end{enumerate} 

    \noindent Now we define how $\tilbordstr_{d,V}$ acts on morphisms:
    \begin{itemize}
    \item[$(\man)$] The morphism $(f,\mathsf{id}_{R})\maps (M,R) \longra (N,R)$ maps the closed submanifold 
    \begin{equation}
    W\subseteq M\times \core(R) \times V
    \end{equation}
    to the pullback along $(f\times \id)$:
    \begin{equation}
        \begin{tikzcd}[ampersand replacement=\&]
            {(f\times \id)^\ast W} \& {W} \\
            \\
            N\times \core(R) \times V \& M\times \core(R) \times V
            \arrow[from=1-1, to=1-2]
            \arrow[hook, "(f\times \id)^\ast \pi_{|W}"',from=1-1, to=3-1]
            \arrow["\lrcorner"{anchor=center, pos=0.125}, draw=none, from=1-1, to=3-2]
            \arrow[hook, from=1-2, to=3-2]
            \arrow["f\times \id"', from=3-1, to=3-2]
        \end{tikzcd}
    \end{equation}
    which is a conically smooth stratified $\big(d + \mathsf{dim}(N)\big)$-dimensional submanifold $f^\ast W \subseteq N \times \core(R) \times V$.
    \item[$(\cR\cube)$] A morphism $(\id_M, R \longra S)$ should act by \emph{forgetting} connected components and restricting the stratified $d$-manifold with corners $(W\ra P)$.
    \end{itemize} 
\end{construction}

\begin{warning}
\Cref{fo:: def: higher strat bord} should be read as a first-pass candidate rather than as a completed definition. The recursive slicing condition captures the geometry one wants, but in its present form it does not yet control where the induced lower-dimensional cuts land after iterating the construction. In particular, this does not yet produce a true functor on the nose.
\end{warning}


The point of \Cref{fo:: def: higher strat bord} is that, if a correct definition of this flavour exists, one expects the $\cube_n$-Segal condition to reduce to \emph{cuts} in a single direction. This is the reason for keeping the construction inductive, and it is precisely this one-directional cutting picture that already appears in the proof of the $(\oo,1)$-case of \Cref{ch5:: thm: stratbord interval segal}.

As a way of wrapping up this section, we will try to motivate our inductive definition by way of examples in $d=2$. Let us start with $R=I() \in \cube_2$ — which corepresents objects in the $\cube_d$ model for $(\oo,d)$-categories. An element of $\tilbordstr_{[0,1,2];V}(\AA^{0},R)$ is a stratified 2-dimensional manifold with corners whose projection to $\core(R) \simeq [0,1]^2$ is a conically smooth submersion as follows:
\medskip
\begin{equation}
    \tikzfig{figures/FO/REC15}{\begin{tikzpicture}
	\begin{pgfonlayer}{nodelayer}
		\node [style=endpoint] (40) at (-5, -1) {};
		\node [style=none] (48) at (-5, -1) {};
		\node [style=none] (4) at (-5, -1) {};
		\node [style=none] (49) at (-3.25, -1) {};
		\node [style=none] (50) at (-10, -6) {};
		\node [style=none] (51) at (-8.25, -6) {};
		\node [style=none] (5) at (3.75, -1) {};
		\node [style=none] (6) at (-3.25, -1) {};
		\node [style=none] (7) at (2, -1) {};
		\node [style=none] (8) at (-10, -6) {};
		\node [style=none] (9) at (-1.25, -6) {};
		\node [style=none] (10) at (-8.25, -6) {};
		\node [style=none] (11) at (-3, -6) {};
		\node [style=endpoint] (41) at (-10, -6) {};
		\node [style=endpoint] (42) at (-1.25, -6) {};
		\node [style=endpoint] (43) at (3.75, -1) {};
		\node [style=none] (44) at (2, -1) {};
		\node [style=none] (45) at (3.75, -1) {};
		\node [style=none] (52) at (2, -1) {};
		\node [style=none] (53) at (3.75, -1) {};
		\node [style=none] (54) at (-3, -6) {};
		\node [style=none] (55) at (-1.25, -6) {};
		\node [style=none] (56) at (-10, -6) {};
		\node [style=none] (59) at (-7, -6) {};
		\node [style=none] (60) at (-4, -6) {};
		\node [style=none] (81) at (-8.25, -6) {};
		\node [style=none] (82) at (-3, -6) {};
		\node [style=none] (83) at (-4.5, -2.25) {};
		\node [style=none] (84) at (2, -1) {};
		\node [style=none] (85) at (-8.25, 4) {};
		\node [style=none] (86) at (-3, 4) {};
		\node [style=none] (87) at (-3.25, -1) {};
		\node [style=none] (88) at (2, -1) {};
		\node [style=none] (89) at (-3.25, 7.75) {};
		\node [style=none] (90) at (2, 7.75) {};
		\node [style=none] (92) at (-8.25, -4.5) {};
		\node [style=none] (114) at (-3.5, -3.25) {};
		\node [style=none] (116) at (-8.25, 0.5) {};
		\node [style=none] (117) at (-3, 0.5) {};
		\node [style=none] (122) at (2, 3.5) {};
		\node [style=none] (123) at (-3.25, 3.25) {};
		\node [style=none] (115) at (-8.25, 2.5) {};
		\node [style=none] (118) at (-5.75, 1.75) {};
		\node [style=none] (119) at (-3, 2.25) {};
		\node [style=none] (120) at (2, 5.25) {};
		\node [style=none] (121) at (-3.25, 6.5) {};
	\end{pgfonlayer}
	\begin{pgfonlayer}{edgelayer}
		\draw [style=ts-dashed-grey] (89.center) to (87.center);
		\draw [style=geometric] (4.center) to (5.center);
		\draw [style={cube_1}] (51.center)
			 to (49.center)
			 to (48.center)
			 to (50.center)
			 to cycle;
		\draw [style={open_int}] (4.center) to (6.center);
		\draw [style={open_int}] (7.center) to (5.center);
		\draw [style=geometric] (8.center) to (9.center);
		\draw [style={open_int}] (8.center) to (10.center);
		\draw [style={open_int}] (11.center) to (9.center);
		\draw [style=red shade low op] (44.center) to (45.center);
		\draw [style={cube_1}] (55.center)
			 to (53.center)
			 to (52.center)
			 to (54.center)
			 to cycle;
		\draw [style=ts-dashed-grey] (81.center) to (82.center);
		\draw [style=ts-dashed-grey] (84.center) to (82.center);
		\draw [style=ts-dashed-grey] (83.center) to (81.center);
		\draw [style=ts-dashed-grey] (85.center) to (81.center);
		\draw [style=ts-dashed-grey] (85.center) to (86.center);
		\draw [style=ts-dashed-grey] (86.center) to (82.center);
		\draw [style=ts-dashed-grey] (89.center) to (90.center);
		\draw [style=ts-dashed-grey] (90.center) to (88.center);
		\draw [style=ts-dashed-grey] (85.center) to (89.center);
		\draw [style=ts-dashed-grey] (86.center) to (90.center);
		\draw [style=geometric shaded] (122.center)
			 to (123.center)
			 to [in=45, out=-150] (116.center)
			 to [in=-165, out=0] (117.center)
			 to cycle;
		\draw [style=geometric shaded] (120.center)
			 to [in=15, out=165] (121.center)
			 to (115.center)
			 to [bend right=15] (118.center)
			 to [bend right=15, looseness=0.75] (119.center)
			 to cycle;
	\end{pgfonlayer}
\end{tikzpicture}
}
\end{equation}
\medskip

\nid and it corresponds to the object which is a disjoint union of a pair of unoriented points. If we consider the case of 1-morphisms which are co-represented by objects of the type $R = I(J_1)$ where $J_1 \in \cube_1 = \Int$ such that $\gamma(J_1) = [0]$. An element of $\tilbordstr_{[0,1,2];V}(\AA^0,R)$ is a stratified 2-dimensional manifold with corners whose projection would look as follows:
\begin{equation}
    \tikzfig{figures/FO/REC16}{\begin{tikzpicture}
	\begin{pgfonlayer}{nodelayer}
		\node [style=0 strat] (108) at (-3.55, 4.125) {};
		\node [style=none] (97) at (-7, 1.75) {};
		\node [style=none] (102) at (-4.75, 2.5) {};
		\node [style=none] (103) at (-2.75, 4.75) {};
		\node [style=none] (104) at (-4.5, 4) {};
		\node [style=endpoint] (40) at (-5, -1) {};
		\node [style=none] (48) at (-5, -1) {};
		\node [style=none] (4) at (-5, -1) {};
		\node [style=none] (49) at (-3.25, -1) {};
		\node [style=none] (50) at (-10, -6) {};
		\node [style=none] (51) at (-8.25, -6) {};
		\node [style=none] (5) at (3.75, -1) {};
		\node [style=none] (6) at (-3.25, -1) {};
		\node [style=none] (7) at (2, -1) {};
		\node [style=none] (8) at (-10, -6) {};
		\node [style=none] (9) at (-1.25, -6) {};
		\node [style=none] (10) at (-8.25, -6) {};
		\node [style=none] (11) at (-3, -6) {};
		\node [style=endpoint] (41) at (-10, -6) {};
		\node [style=endpoint] (42) at (-1.25, -6) {};
		\node [style=endpoint] (43) at (3.75, -1) {};
		\node [style=none] (44) at (2, -1) {};
		\node [style=none] (45) at (3.75, -1) {};
		\node [style=none] (52) at (2, -1) {};
		\node [style=none] (53) at (3.75, -1) {};
		\node [style=none] (54) at (-3, -6) {};
		\node [style=none] (55) at (-1.25, -6) {};
		\node [style=none] (56) at (-10, -6) {};
		\node [style=none] (57) at (-5.75, -4.75) {};
		\node [style=none] (58) at (-2.75, -4.75) {};
		\node [style=none] (59) at (-7, -6) {};
		\node [style=none] (60) at (-4, -6) {};
		\node [style=none] (69) at (-2, -1) {};
		\node [style=none] (70) at (0.75, -1) {};
		\node [style=none] (71) at (-3.25, -2.25) {};
		\node [style=none] (72) at (-0.5, -2.25) {};
		\node [style=none] (81) at (-7, -4.75) {};
		\node [style=none] (82) at (-1.75, -4.75) {};
		\node [style=none] (83) at (-4.5, -2.25) {};
		\node [style=none] (84) at (0.75, -2.25) {};
		\node [style=none] (85) at (-7, 5.25) {};
		\node [style=none] (86) at (-1.75, 5.25) {};
		\node [style=none] (87) at (-4.5, -2.25) {};
		\node [style=none] (88) at (0.75, -2.25) {};
		\node [style=none] (89) at (-4.5, 7.75) {};
		\node [style=none] (90) at (0.75, 7.75) {};
		\node [style=none] (92) at (-7, -3.25) {};
		\node [style=none] (95) at (-4.5, -1) {};
		\node [style=none] (96) at (-7, 3.5) {};
		\node [style=none] (98) at (-4.5, 5.5) {};
		\node [style=none] (100) at (-4.75, 2.5) {};
		\node [style=none] (101) at (-2.75, 4.75) {};
		\node [style=0 strat] (105) at (-5.25, 3) {};
		\node [style=0 strat] (106) at (-3.25, 5.25) {};
		\node [style=0 strat] (107) at (-5.5, 2) {};
		\node [style=none] (109) at (-4.95, 2.8) {};
		\node [style=none] (110) at (-1.75, 1) {};
		\node [style=none] (111) at (-1.75, 2.5) {};
		\node [style=none] (112) at (-3.75, 2.25) {};
		\node [style=none] (113) at (-1.5, 4.25) {};
		\node [style=none] (114) at (0.75, 3) {};
		\node [style=none] (115) at (-1.75, 2.5) {};
		\node [style=none] (116) at (0.75, 4.25) {};
		\node [style=none] (117) at (-1.5, 4.25) {};
		\node [style=none] (118) at (-3.75, 2.25) {};
	\end{pgfonlayer}
	\begin{pgfonlayer}{edgelayer}
		\draw [style=geometric, in=225, out=30] (109.center) to (108);
		\draw [style=geometric shaded] (104.center)
			 to (97.center)
			 to [bend right=60, looseness=0.50] (102.center)
			 to (103.center)
			 to [bend left] cycle;
		\draw [style=ts-dashed-grey] (89.center) to (87.center);
		\draw [style=geometric] (4.center) to (5.center);
		\draw [style={cube_1}] (51.center)
			 to (49.center)
			 to (48.center)
			 to (50.center)
			 to cycle;
		\draw [style={open_int}] (4.center) to (6.center);
		\draw [style={open_int}] (7.center) to (5.center);
		\draw [style=geometric] (8.center) to (9.center);
		\draw [style={open_int}] (8.center) to (10.center);
		\draw [style={open_int}] (11.center) to (9.center);
		\draw [style=red shade low op] (44.center) to (45.center);
		\draw [style={cube_1}] (55.center)
			 to (53.center)
			 to (52.center)
			 to (54.center)
			 to cycle;
		\draw [style={cube_1}] (60.center)
			 to (58.center)
			 to (57.center)
			 to (59.center)
			 to cycle;
		\draw [style={cube_1}] (72.center)
			 to (70.center)
			 to (69.center)
			 to (71.center)
			 to cycle;
		\draw [style={open_int}] (59.center) to (60.center);
		\draw [style={open_int}] (69.center) to (70.center);
		\draw [style=ts-dashed-grey] (81.center) to (82.center);
		\draw [style=ts-dashed-grey] (83.center) to (84.center);
		\draw [style=ts-dashed-grey] (84.center) to (82.center);
		\draw [style=ts-dashed-grey] (83.center) to (81.center);
		\draw [style=ts-dashed-grey] (85.center) to (81.center);
		\draw [style=ts-dashed-grey] (85.center) to (86.center);
		\draw [style=ts-dashed-grey] (86.center) to (82.center);
		\draw [style=ts-dashed-grey] (87.center) to (88.center);
		\draw [style=ts-dashed-grey] (89.center) to (90.center);
		\draw [style=ts-dashed-grey] (90.center) to (88.center);
		\draw [style=ts-dashed-grey] (85.center) to (89.center);
		\draw [style=ts-dashed-grey] (86.center) to (90.center);
		\draw [style=geometric shaded] (98.center)
			 to (96.center)
			 to [bend left=45, looseness=0.75] (100.center)
			 to (101.center)
			 to [bend right=60, looseness=0.75] cycle;
		\draw [style=geometric, in=225, out=60] (105) to (106);
		\draw [style=geometric, in=-150, out=60, looseness=0.75] (107) to (109.center);
		\draw [style=geometric shaded] (113.center)
			 to [in=150, out=-120, looseness=1.50] (114.center)
			 to (110.center)
			 to [in=-90, out=150, looseness=1.25] (112.center)
			 to cycle;
		\draw [style=geometric shaded] (118.center)
			 to [in=180, out=45] (115.center)
			 to (116.center)
			 to [in=45, out=180, looseness=0.50] (117.center)
			 to cycle;
	\end{pgfonlayer}
\end{tikzpicture}
}
\end{equation}
\medskip

\nid Finally, to represent a 2-morphism, one should allow a cube $R = I(J_1)\in \cube_2$ where $J_1 $ has $\gamma(J_1)=[1]$ and whose two first-direction components carry interval data in the transverse direction. An element of $\tilbordstr_{[0,1,2];V}(\AA^0,R)$ should then be thought of as a stratified surface interpolating between 1-morphisms of the previous kind, as in the following picture:

\begin{equation}
    \tikzfig{figures/FO/REC17}{\begin{tikzpicture}
	\begin{pgfonlayer}{nodelayer}
		\node [style=none] (107) at (-7, 2.25) {};
		\node [style=none] (108) at (-4.5, 5) {};
		\node [style=none] (109) at (-5.25, 2.75) {};
		\node [style=none] (110) at (0.5, 5) {};
		\node [style=none] (111) at (-2, 2.25) {};
		\node [style=none] (112) at (-4.5, 2.75) {};
		\node [style=none] (113) at (-3.5, 4) {};
		\node [style=endpoint] (40) at (-5, -1) {};
		\node [style=none] (48) at (-5, -1) {};
		\node [style=none] (4) at (-5, -1) {};
		\node [style=none] (49) at (-3.25, -1) {};
		\node [style=none] (50) at (-10, -6) {};
		\node [style=none] (51) at (-8.25, -6) {};
		\node [style=none] (5) at (3.75, -1) {};
		\node [style=none] (6) at (-3.25, -1) {};
		\node [style=none] (7) at (2, -1) {};
		\node [style=none] (8) at (-10, -6) {};
		\node [style=none] (9) at (-1.25, -6) {};
		\node [style=none] (10) at (-8.25, -6) {};
		\node [style=none] (11) at (-3, -6) {};
		\node [style=endpoint] (41) at (-10, -6) {};
		\node [style=endpoint] (42) at (-1.25, -6) {};
		\node [style=endpoint] (43) at (3.75, -1) {};
		\node [style=none] (44) at (2, -1) {};
		\node [style=none] (45) at (3.75, -1) {};
		\node [style=none] (52) at (2, -1) {};
		\node [style=none] (53) at (3.75, -1) {};
		\node [style=none] (54) at (-3, -6) {};
		\node [style=none] (55) at (-1.25, -6) {};
		\node [style=none] (56) at (-10, -6) {};
		\node [style=none] (57) at (-5.75, -4.75) {};
		\node [style=none] (58) at (-2.75, -4.75) {};
		\node [style=none] (59) at (-7, -6) {};
		\node [style=none] (60) at (-4, -6) {};
		\node [style=none] (69) at (-2, -1) {};
		\node [style=none] (70) at (0.75, -1) {};
		\node [style=none] (71) at (-3.25, -2.25) {};
		\node [style=none] (72) at (-0.5, -2.25) {};
		\node [style=none] (77) at (-3.75, -2.75) {};
		\node [style=none] (78) at (-1, -2.75) {};
		\node [style=none] (79) at (-5, -4) {};
		\node [style=none] (80) at (-2.25, -4) {};
		\node [style=none] (81) at (-7, -4.75) {};
		\node [style=none] (82) at (-1.75, -4.75) {};
		\node [style=none] (83) at (-4.5, -2.25) {};
		\node [style=none] (84) at (0.75, -2.25) {};
		\node [style=none] (85) at (-7, 5.25) {};
		\node [style=none] (86) at (-1.75, 5.25) {};
		\node [style=none] (87) at (-4.5, -2.25) {};
		\node [style=none] (88) at (0.75, -2.25) {};
		\node [style=none] (89) at (-4.5, 7.75) {};
		\node [style=none] (90) at (0.75, 7.75) {};
		\node [style=none] (92) at (-7, -3.25) {};
		\node [style=none] (95) at (-4.5, -1) {};
		\node [style=none] (100) at (-7, 3.75) {};
		\node [style=none] (101) at (-5.25, 2.75) {};
		\node [style=none] (102) at (-3.5, 4) {};
		\node [style=none] (103) at (-4.5, 2.75) {};
		\node [style=none] (104) at (-1.75, 3.5) {};
		\node [style=none] (105) at (0.5, 6.25) {};
		\node [style=none] (106) at (-4.5, 6.25) {};
		\node [style=none] (114) at (-3.5, -3.25) {};
		\node [style=0 strat] (115) at (-5.725, 3.35) {};
		\node [style=0 strat] (116) at (-5.75, 2.325) {};
		\node [style=0 strat] (118) at (-3.25, 4.75) {};
		\node [style=0 strat] (119) at (-1.75, 6.25) {};
		\node [style=0 strat] (122) at (-4.75, 4.5) {};
		\node [style=none] (123) at (-4.75, 3.375) {};
	\end{pgfonlayer}
	\begin{pgfonlayer}{edgelayer}
		\draw [style=ts-dashed-grey] (89.center) to (87.center);
		\draw [style=geometric shaded] (108.center)
			 to (110.center)
			 to (111.center)
			 to [in=-120, out=-180, looseness=0.75] (112.center)
			 to [in=-15, out=60] (113.center)
			 to [in=45, out=165] (109.center)
			 to [in=0, out=-90, looseness=0.50] (107.center)
			 to cycle;
		\draw [style=box, in=60, out=180] (123.center) to (116);
		\draw [style=ts-dashed-red] (102.center) to (114.center);
		\draw [style=geometric] (4.center) to (5.center);
		\draw [style={cube_1}] (51.center)
			 to (49.center)
			 to (48.center)
			 to (50.center)
			 to cycle;
		\draw [style={open_int}] (4.center) to (6.center);
		\draw [style={open_int}] (7.center) to (5.center);
		\draw [style=geometric] (8.center) to (9.center);
		\draw [style={open_int}] (8.center) to (10.center);
		\draw [style={open_int}] (11.center) to (9.center);
		\draw [style=red shade low op] (44.center) to (45.center);
		\draw [style={cube_1}] (55.center)
			 to (53.center)
			 to (52.center)
			 to (54.center)
			 to cycle;
		\draw [style={cube_1}] (60.center)
			 to (58.center)
			 to (57.center)
			 to (59.center)
			 to cycle;
		\draw [style={cube_1}] (72.center)
			 to (70.center)
			 to (69.center)
			 to (71.center)
			 to cycle;
		\draw [style={cube_1}] (80.center)
			 to (78.center)
			 to (77.center)
			 to (79.center)
			 to cycle;
		\draw [style={open_int}] (59.center) to (60.center);
		\draw [style={open_int}] (69.center) to (70.center);
		\draw [style=ts-dashed-grey] (81.center) to (82.center);
		\draw [style=ts-dashed-grey] (83.center) to (84.center);
		\draw [style=ts-dashed-grey] (84.center) to (82.center);
		\draw [style=ts-dashed-grey] (83.center) to (81.center);
		\draw [style=ts-dashed-grey] (85.center) to (81.center);
		\draw [style=ts-dashed-grey] (85.center) to (86.center);
		\draw [style=ts-dashed-grey] (86.center) to (82.center);
		\draw [style=ts-dashed-grey] (87.center) to (88.center);
		\draw [style=ts-dashed-grey] (89.center) to (90.center);
		\draw [style=ts-dashed-grey] (90.center) to (88.center);
		\draw [style=ts-dashed-grey] (85.center) to (89.center);
		\draw [style=ts-dashed-grey] (86.center) to (90.center);
		\draw [style=geometric shaded] (102.center)
			 to [in=45, out=165] (101.center)
			 to [in=-15, out=90, looseness=0.75] (100.center)
			 to (106.center)
			 to (105.center)
			 to (104.center)
			 to [in=60, out=180, looseness=0.75] (103.center)
			 to [in=-15, out=60] cycle;
		\draw [style=box, in=30, out=-30, looseness=0.75] (122) to (123.center);
		\draw [style=box, in=180, out=45] (122) to (118);
		\draw [style=box] (115) to (122);
		\draw [style=box, in=-150, out=30, looseness=1.25] (118) to (119);
	\end{pgfonlayer}
\end{tikzpicture}
}
\end{equation}
\bigskip

\section{Stratified TQFTs via factorization homology}

In the unstratified setting, factorization homology is one of the cleanest mechanisms we have for passing from local algebraic data to topological field theories. The reason is simple: every smooth $d$-manifold is locally modelled on $\RR^d$. One may therefore encode local observables by a single algebraic object --- an $\cE_d$-algebra, or equivalently a $\Disk_d$-algebra --- and then integrate that datum over the manifold. In the fully extended setting, this perspective is already present in Lurie's sketch and is later developed in Scheimbauer's thesis \cite{Lur09,LurieHA,AF1,Sch14}.

For stratified spaces, however, there is no longer a unique local model. Near a point in a codimension-$k$ stratum one sees a basic of the form $\RR^{d-k}\times \sC(L)$, and the link $L$ remembers how the deeper strata sit nearby. So the coefficient system must now keep track not only of the top-dimensional disk, but of all the allowed singularity types and the embeddings among them. This is the conceptual reason for replacing ordinary disk algebras by $\cB$-disk algebras for a suitable category of basics $\cB$.

There is also a second language lurking in the background, which is often better adapted to field-theoretic intuition. Rather than presenting the coefficient system directly as a disk-algebra, one can recast the same local-to-global structure in the language of factorization algebras on opens. In the stratified setting the relevant objects should be \emph{constructible factorization algebras}: roughly, factorization algebras that are locally constant along each stratum and whose variation across strata is controlled by the stratification. We will not try to develop that formalism here, but it is worth keeping it in mind, since it is the natural language for pushing observables along the projection of a stratified bordism to an interval. With that motivation in place, let us recall the operadic definition from \cite{AFT3}.

\begin{definition}\cite[Definition~2.2 and Notation~2.9]{AFT3}
    Let $\cB \longra \Bsc$ be an $(\oo,1)$-category of basics and let $\cV^{\otimes}$ be a symmetric monoidal $(\oo,1)$-category. The symmetric monoidal $(\oo,1)$-category $\Disk(\cB)\subset \Mfd(\cB)$ is the smallest full symmetric monoidal subcategory generated by the basics in $\cB$. A \define{$\cB$-disk algebra in} $\cV$ is a symmetric monoidal functor
    \begin{equation}
        A \maps \Disk(\cB) \longra \cV.
    \end{equation}
\end{definition}

\begin{definition}\cite[Definition~2.14 and Theorem~2.15]{AFT3}
    Let $A$ be a $\cB$-disk algebra in a symmetric monoidal $(\oo,1)$-category $\cV^{\otimes}$ for which factorization homology exists, and let $X\in \Mfd(\cB)$. The \define{factorization homology of $X$ with coefficients in $A$} is the left Kan extension of $A$ along $\Disk(\cB)\hookra \Mfd(\cB)$, and is denoted
    \begin{equation}
        \int_{X} A.
    \end{equation}
    When $\cV^{\otimes}$ is $\otimes$-sifted cocomplete, it is computed by the colimit formula
    \begin{equation}
        \int_{X} A \,\simeq\, \colim_{(U\hookra X)\in \Disk(\cB)_{/X}} A(U).
    \end{equation}
\end{definition}

The picture to keep in mind is the following. Once a defect datum
\begin{equation}
    \Bsc_{d}^{\cS_{d-1},\DD} \longra \Bsc_{d}^{\cS_{d-1}}
\end{equation}
has been fixed, the relevant local coefficients should be $\Bsc_{d}^{\cS_{d-1},\DD}$-disk algebras. In other words, one does not attach algebraic data only to the smooth $d$-disk, but to every allowed local singularity type, together with the compatibilities imposed by embeddings of basics. This is precisely the point where the defect labels introduced in Chapter~5 begin to interact with factorization homology in a meaningful way.

Now let $W$ be a morphism in our interval Segal space $\bordstr_d$, together with its projection $\pi\maps W \longra \core(I)$. If $A$ is such a stratified disk algebra, one expects the assignment
\begin{equation}
    U \subset \core(I) \longmapsto \int_{\pi^{-1}(U)} A
\end{equation}
to define a constructible factorization algebra on $\core(I)$ relative to the finite stratification induced by $I$. If $I\in \Int$ encodes a finite family of cuts, let $S_{I}\subset \core(I)$ denote the corresponding finite set of separating points. Then the expected pushforward should be locally constant on each connected component of $\core(I)\setminus S_{I}$ and only allowed to jump across the strata in the induced stratification
\begin{equation}
    \emptyset \subset S_{I} \subset \core(I).
\end{equation}
This is precisely the kind of one-dimensional constructibility one wants in order to remember the algebraic data assigned to composition. In the unstratified case, Scheimbauer uses factorization algebras on $(0,1)$ which are locally constant with respect to such finite stratifications to model bimodules and, more generally, the levels of the Morita target \cite[Corollary~3.2.25 and Definition~3.2.27]{Sch14}. The same pattern suggests that in the stratified setting the first target should not yet be $\Vect$, but rather an interval Segal space of constructible factorization algebras on stratified intervals. Moreover, if one cuts $W$ along a regular value of $\pi$, the excision theorem of \cite{AFT3} identifies the factorization homology of the whole bordism with the relative tensor product of the factorization homologies of the two pieces. Thus the gluing law required for composition is not added by hand; it is already built into the formalism.

\begin{conjecture}
    Let $\cV^{\otimes}$ be a symmetric monoidal $(\oo,1)$-category with the colimits needed to define stratified factorization homology, and let $A$ be a $\Bsc_{d}^{\cS_{d-1},\DD}$-disk algebra in $\cV$. Then factorization homology with coefficients in $A$ should determine a symmetric monoidal functor out of the stratified bordism category constructed in Chapter~5, landing naturally in an interval Segal space of constructible factorization algebras on stratified intervals. After comparison with the corresponding Morita-type target, this should produce a non-extended stratified topological quantum field theory.
\end{conjecture}

















